\documentclass[11pt]{article}

\usepackage[margin=1in]{geometry}
\usepackage{amsmath,amssymb,amsthm}
\usepackage{booktabs,microtype}
\usepackage{xcolor}
\usepackage{url}
\usepackage[colorlinks=true,linkcolor=blue!55!black,
  citecolor=green!45!black,urlcolor=blue!55!black]{hyperref}
\usepackage[nameinlink,noabbrev]{cleveref}
\numberwithin{equation}{section}

\newtheorem{theorem}{Theorem}[section]
\newtheorem{lemma}[theorem]{Lemma}
\newtheorem{proposition}[theorem]{Proposition}
\newtheorem{corollary}[theorem]{Corollary}
\newtheorem{problem}[theorem]{Problem}
\newtheorem{maintheorem}{Theorem}
\renewcommand{\themaintheorem}{\Alph{maintheorem}}
\theoremstyle{definition}
\newtheorem{definition}[theorem]{Definition}
\newtheorem{example}[theorem]{Example}
\theoremstyle{remark}
\newtheorem{remark}[theorem]{Remark}

\newcommand{\Damp}{\mathcal D_{\mathrm{amp}}}
\newcommand{\Sh}{\operatorname{Sh}}
\newcommand{\SSh}{\operatorname{SSh}}
\newcommand{\cmin}{c_{\min}}
\newcommand{\mamp}{m_{\mathrm{amp}}}
\newcommand{\unsat}{\textsc{unsat}}

\title{Corners of Small Ample Classes:\\ Facet Cubes and Certified Searches}
\author{Olivier Bousquet}
\date{}

\hypersetup{
  pdftitle={Corners of Small Ample Classes: Facet Cubes and Certified Searches},
  pdfauthor={Olivier Bousquet}
}

\begin{document}
\maketitle

\begin{abstract}
A vertex of an ample class is a corner if it lies in a unique maximal
cube.  Hall's maximum class of VC dimension three, with \(299\) concepts on
twelve coordinates, has no corner, so ample classes need not admit corner
peelings.  We ask how small a nonempty ample class without corners can be.
This least order equals the least order of an ample forbidden
partial-cube minor for corner peelability, which was known to lie between
\(64\) and \(291\).  We first observe that an ample class is the union of
exactly one cube for each maximal shattered set, that these are its maximal
cubes, and that its corners are the vertices covered once.  Conversely,
choosing one cube for each facet of a simplicial complex produces an ample
class whenever the union shatters nothing outside the complex.  This gives
an exact propositional description of cornerless ample classes, which we
combine with lexicographic symmetry breaking and DRAT proof checking.  We
prove that every maximum class on at most eight coordinates has a corner,
that every nonempty ample class on at most seven coordinates has a corner,
and that every cornerless ample one-coordinate lift of a single or double
contraction of Hall's class has at least \(291\)
concepts.  For single contractions we determine the optimum exactly: it is
\(291\) for eight of the twelve coordinates and \(295\) for the other four.
We strengthen the seven-coordinate result by proving that every such ample
class can be corner-peeled to any prescribed vertex. The proof excludes
unique corners through an exhaustive, certified split of their fibers.
We also record that Hall's class has no nontrivial cube symmetry, and that
its twelve coordinates fall into six pairs with equal contractions, a
pairing also singled out by the corner counts of double contractions.
Finally, iterating contraction and optimal lifting from these classes
produces a cornerless ample class of order \(267\) on twelve coordinates,
which lowers the upper bound to \(\cmin\le267\).
\end{abstract}

\medskip
\noindent\textbf{Keywords.}
Ample class, lopsided set, maximum class, corner peeling, VC dimension,
partial cube, SAT solving, DRAT proofs.

\tableofcontents

%======================================================================
\section{Introduction}
\label{sec:intro}

\subsection{Ample classes and corners}

Let \(E\) be a finite set of \emph{coordinates} and let
\(Q_E=\{0,1\}^E\) be the hypercube on \(E\).  A \emph{class} (or
\emph{family}) is a subset \(C\subseteq Q_E\); its elements are called
\emph{concepts} or \emph{vertices}, and \(|C|\) is its \emph{order}.  A class
\emph{shatters} \(S\subseteq E\) if its projection to \(S\) is all of
\(Q_S\), and it \emph{strongly shatters} \(S\) if it contains a whole cube of
dimension \(|S|\) in which exactly the coordinates in \(S\) vary.  A class is
\emph{ample} if these two notions agree.  Ample classes were introduced by
Lawrence as \emph{lopsided sets} \cite{Lawrence83} and were later studied
under this name and as extremal or shattering-extremal classes
\cite{BandeltCDK06,Bollobas95,Anstee2002,Moran2012}.  They include the
\emph{maximum} classes, those meeting the Sauer--Shelah bound with
equality \cite{Sauer1972Density,Floyd95}.

A vertex of \(C\) is a \emph{corner} if it lies in a unique maximal cube of
\(C\).  Deleting a corner from an ample class leaves an ample class
\cite[Lemma~4.1]{ChalopiChepoiMoran22}.  A \emph{corner peeling} deletes
corners one at a time until nothing is left.  Corner peelings of maximum
classes were proposed as a route to unlabeled sample compression schemes
\cite{Kuzmin06,Rubinstein2009a}.  Chalopin, Chepoi, Moran and Warmuth showed
that a maximum class of VC dimension three with \(299\) concepts on twelve
coordinates, derived from Hall's nonextendable partial shelling of the
twelve-dimensional cross-polytope \cite{Hall04}, has no corner at all
\cite[Theorem~4.5]{ChalopiChepoiMoran22}.  We call it \emph{Hall's class}
and write \(C_H\).  They also proved that ample classes of VC dimension at
most two always admit corner peelings
\cite[Proposition~4.11]{ChalopiChepoiMoran22}.

\subsection{The question}

Hall's class shows that cornerless ample classes exist.  It leaves open how
small they can be.

\begin{definition}
\label{def:cmin}
Let \(\cmin\) be the least order of a nonempty ample class with no corner.
\end{definition}

This number has a second meaning.  Ample classes are identified with
partial cubes, and the class \(\Damp\) of corner-peelable ample classes is
closed under taking partial-cube minors (pc-minors)
\cite{KnauerMarc20,BousquetDampReader}.  Let \(\mamp\) be the least order of
an ample forbidden pc-minor of \(\Damp\).
Proposition~\ref{prop:cmin-equals-mamp} shows that
\[
 \cmin=\mamp .
\]
The companion work \cite{BousquetDampReader} proves
\(64\le\mamp\le291\); the upper bound comes from an explicit family, and
the lower bound combines structural reductions with finite exclusions.
This paper lowers the upper bound to \(267\) (Theorem~\ref{thm:D}) and rules
out cornerless ample classes on few coordinates and near Hall's class.

\subsection{Results}

Our main tool is elementary.  For an ample class, the maximal cubes
correspond bijectively to the maximal shattered sets, and the class is the
union of these cubes (Lemma~\ref{lem:facet-cubes}).  Thus a corner is a
vertex covered by exactly one of them, and a cornerless ample class is a
choice of cubes that covers every vertex of its union at least twice.  The
converse half of the lemma turns this into an exact finite search: a
choice of one cube per facet of a simplicial complex \(\Sigma\) is ample,
with shattered sets exactly \(\Sigma\), as soon as the union shatters no
set outside \(\Sigma\).  No ampleness test appears in the search; the
Sandwich inequalities supply it.

We encode this description in propositional logic, break the
hyperoctahedral symmetry with lexicographic-leader constraints, and check
every unsatisfiability claim with an independently verified DRAT proof.

\begin{maintheorem}[Maximum classes]
\label{thm:A}
Every maximum class on at most eight coordinates has a corner.
\end{maintheorem}

\begin{maintheorem}[Few coordinates]
\label{thm:B}
Every nonempty ample class on at most seven coordinates has a corner.
Equivalently, every nonempty cornerless ample class has at least eight
coordinates on which it is not constant.
\end{maintheorem}

\begin{maintheorem}[Lifts of Hall contractions]
\label{thm:C}
Let \(B\) be the contraction of Hall's class at one coordinate or at two
coordinates.  If \(Y\subseteq B\times\{0,1\}\) is a cornerless ample class
whose projection to the first factor is \(B\), then \(|Y|\ge291\).  For
the contraction \(C_H/e\) at one coordinate \(e\in\{0,\dots,11\}\), the least
order of such a lift is exactly
\[
 291\ \text{ if } e\le7,\qquad 295\ \text{ if } e\ge8 .
\]
\end{maintheorem}

\begin{maintheorem}[A smaller cornerless class]
\label{thm:D}
There is a cornerless ample class of VC dimension \(3\) and order \(267\) on
twelve coordinates.  Consequently
\[
 64\le\cmin=\mamp\le267 .
\]
\end{maintheorem}

\begin{maintheorem}[Prescribed endpoints]
\label{thm:E}
Every nonempty ample class on at most seven coordinates can be
corner-peeled to any prescribed vertex. In particular, every nontrivial
such class has at least two corners.
\end{maintheorem}

Theorems~\ref{thm:A}--\ref{thm:C} and~\ref{thm:E} are computer-assisted.  The mathematical
reductions from each statement to finite formulas are proved in full; what
is delegated to the computer is the unsatisfiability of explicit formulas,
and every such claim is backed by a checked proof
(Section~\ref{sec:certification}).  Theorem~\ref{thm:C} is sharp for
single contractions and recovers the upper bound \(\mamp\le291\) by an
independent route: a satisfying assignment at bound \(291\) is itself a
cornerless ample class (Proposition~\ref{prop:291}).

\begin{table}[t]
\centering
\small
\begin{tabular}{@{}p{0.29\textwidth}p{0.33\textwidth}p{0.28\textwidth}@{}}
\toprule
Statement & Scope & Method\\
\midrule
Lemma~\ref{lem:facet-cubes} (facet cubes) & every ample class & proof\\
Proposition~\ref{prop:cmin-equals-mamp} (\(\cmin=\mamp\)) & all orders & proof\\
Theorem~\ref{thm:encoding} (exactness) & all \(n,d\) & proof\\
Theorem~\ref{thm:A} & maximum, \(\le8\) coordinates & Theorem~\ref{thm:B} + \(4\) certificates\\
Theorem~\ref{thm:B} & ample, \(\le7\) coordinates & proof + \(15\) certificates\\
Theorem~\ref{thm:C} & lifts of \(78\) Hall contractions & proof + \(80\) certificates\\
Proposition~\ref{prop:291} & orders \(291\), \(295\) & direct check\\
Proposition~\ref{prop:hall-structure} & Hall's class & direct check\\
Theorem~\ref{thm:D} & order \(267\) & direct check\\
Theorem~\ref{thm:E} & prescribed endpoint, \(\le7\) coordinates & proof + fiber certificates\\
\bottomrule
\end{tabular}
\caption{Results at a glance.  A certificate is a DRAT refutation of an
explicit formula, accepted by \texttt{drat-trim}.}
\label{tab:glance}
\end{table}

\subsection{Where the results sit}

Theorem~\ref{thm:B} constrains the number of coordinates, while the lower
bound \(\cmin\ge64\) of \cite{BousquetDampReader} constrains the order;
neither implies the other.  A class on eight coordinates can have up to
\(256\) vertices, so Theorem~\ref{thm:B} alone gives no order bound beyond
the trivial one, and the order bound alone does not exclude seven
coordinates, where classes of order up to \(128\) exist.  Hall's class uses
twelve coordinates.

Theorem~\ref{thm:A} complements the known low-dimensional results.
Maximum classes of VC dimension at most two peel
\cite[Proposition~4.11]{ChalopiChepoiMoran22}, and those of VC dimension
\(n-1\) and \(n\) have corners (Example~\ref{ex:codim-one}).  For VC dimension
three, the companion work proves the case of at most eight coordinates by a
different encoding \cite{BousquetDampReader}; here we recertify that case
with the encoding of Section~\ref{sec:encoding}.

Theorem~\ref{thm:C} explores the neighbourhood of the only known small
cornerless class.  Every cornerless ample class \(Y\) with a coordinate
\(z\) has the form \((P\times\{0\})\cup(Q\times\{1\})\) with
\(P\cup Q\) its contraction at \(z\).  Theorem~\ref{thm:C} says that no
class below order \(291\) arises this way over a single or double
contraction of Hall's class.  The companion work extends the same search to
\(324\) further bases in a larger neighbourhood of Hall's class
\cite{BousquetDampReader}.

\subsection{Organization}

Section~\ref{sec:prelim} fixes notation and gives running examples.
Section~\ref{sec:facet} proves the facet-cube lemma and the identity
\(\cmin=\mamp\).  Section~\ref{sec:encoding} turns the lemma into formulas
and proves that they are exact.  Section~\ref{sec:certification} explains
how unsatisfiability is certified and what is trusted.
Sections~\ref{sec:maximum}--\ref{sec:hall} prove
Theorems~\ref{thm:A}--\ref{thm:C}.  Section~\ref{sec:open} discusses the
limits of the method and open problems.

%======================================================================
\section{Preliminaries}
\label{sec:prelim}

\subsection{Cubes, shattering and ampleness}

We identify \(Q_E\) with the power set of \(E\) when convenient, and write
vertices of \(Q_{[n]}\), \([n]=\{1,\dots,n\}\), as binary words
\(x_1x_2\cdots x_n\).  For \(T\subseteq E\) and a vertex \(b\) with
\(b_i=0\) for all \(i\in T\), the \emph{cube} with \emph{support} \(T\)
and \emph{placement} \(b\) is
\[
 Q(T,b)=\{x\in Q_E: x_i=b_i\text{ for all } i\notin T\}.
\]
It has \(2^{|T|}\) vertices.  A vertex \(v\) lies in the cube with support
\(T\) exactly when the placement is \(v\setminus T\), the vertex obtained
from \(v\) by setting the coordinates in \(T\) to zero.  Thus \(v\) lies in
exactly one cube of each support.

For \(C\subseteq Q_E\), let \(\Sh(C)\) be the family of sets shattered by
\(C\) and \(\SSh(C)\) the family of sets strongly shattered by \(C\), that
is, the supports of cubes contained in \(C\).  Both families are
down-closed, and \(\SSh(C)\subseteq\Sh(C)\).  The \emph{VC dimension} of
\(C\) is the largest size of a member of \(\Sh(C)\).  A coordinate \(e\) is
\emph{active} in \(C\) if \(C\) is not constant on it, that is, if
\(\{e\}\in\Sh(C)\).

\begin{theorem}[Sandwich inequalities {\cite{Bollobas95,Anstee2002,BandeltCDK06}}]
\label{thm:sandwich}
For every \(C\subseteq Q_E\),
\[
 |\SSh(C)|\le|C|\le|\Sh(C)|.
\]
Equality holds in either inequality if and only if \(C\) is ample,
that is, if and only if \(\SSh(C)=\Sh(C)\)
\cite[Corollary~1 and Theorem~2]{BandeltCDK06}.
\end{theorem}

For \(0\le d\le n\), put \(\Phi_d(n)=\sum_{i=0}^d\binom ni\).  A class of VC
dimension \(d\) on \(n\) coordinates has at most \(\Phi_d(n)\) vertices
\cite{Sauer1972Density}; it is \emph{maximum} if equality holds.  By
Theorem~\ref{thm:sandwich}, a class \(C\) of VC dimension \(d\) is maximum
if and only if it is ample and \(\Sh(C)\) consists of all sets of size at
most \(d\): indeed \(\Phi_d(n)=|C|\le|\Sh(C)|\le\Phi_d(n)\).

For \(T\subseteq E\), the \emph{reduction} of \(C\) along \(T\) is the set
of placements of the \(T\)-cubes in \(C\),
\[
 C^T=\{b\in Q_{E\setminus T}: Q(T,b)\subseteq C\},
\]
where we view placements as vertices of \(Q_{E\setminus T}\).  We use two
standard facts about ample classes.

\begin{theorem}[{\cite[Theorems~3.1 and~3.2]{ChalopiChepoiMoran22}}]
\label{thm:reductions}
Let \(C\) be ample and \(T\subseteq E\).  Then \(C^T\) is ample, and if it
is nonempty it is connected in the hypercube graph.
\end{theorem}

The shattering of a reduction is read off directly from that of \(C\): a
cube with support \(U\subseteq E\setminus T\) in \(C^T\) is the same as a
cube with support \(U\cup T\) in \(C\), so
\begin{equation}
 \SSh(C^T)=\{U\subseteq E\setminus T:\ U\cup T\in\SSh(C)\}.
 \label{eq:ssh-reduction}
\end{equation}

\subsection{Corners, peelings and pc-minors}

A cube \(Q(T,b)\subseteq C\) is \emph{maximal} if no cube of \(C\)
properly contains it.  A vertex \(v\in C\) is a \emph{corner} if exactly one
maximal cube of \(C\) contains \(v\).  The class is \emph{cornerless} if it
has no corner.  A \emph{corner peeling} of \(C\) is an ordering
\(v_1,\dots,v_m\) of \(C\) such that \(v_i\) is a corner of
\(C\setminus\{v_1,\dots,v_{i-1}\}\) for every \(i\).

\begin{theorem}[{\cite[Lemma~4.1]{ChalopiChepoiMoran22}}]
\label{thm:corner-deletion}
If \(C\) is ample and \(v\in C\), then \(C\setminus\{v\}\) is ample if and
only if \(v\) is a corner of \(C\).
\end{theorem}

Consequently every class that has a corner peeling is ample, and every
proper prefix deletion stays ample.  We write \(\Damp\) for the class of
nonempty ample classes that admit a corner peeling.

For \(e\in E\), the \emph{contraction} \(\pi_eC\subseteq Q_{E\setminus\{e\}}\)
deletes coordinate \(e\) from every vertex, and the \emph{restriction}
\(\rho_e^iC\) keeps the vertices with \(x_e=i\) and then deletes \(e\).
A \emph{pc-minor} of \(C\) is a class obtained by a finite sequence of
contractions and nonempty restrictions, with inactive coordinates
discarded \cite{KnauerMarc20}; it is \emph{proper} if it is not isomorphic
to \(C\).  Ample classes are closed under both operations
\cite[Theorem~3.1]{ChalopiChepoiMoran22}, and the class \(\Damp\) is
closed under pc-minors \cite[Proposition~2.1]{BousquetDampReader}.  An
\emph{ample forbidden pc-minor} of \(\Damp\) is an ample class not in
\(\Damp\) all of whose proper pc-minors lie in \(\Damp\).

\subsection{Running examples}

\begin{example}[A cube minus a vertex]
\label{ex:codim-one}
Let \(C=Q_3\setminus\{111\}\).  It shatters every set of size at most two
but not \(\{1,2,3\}\), and it contains the three squares
\[
 Q(\{1,2\},000),\qquad Q(\{1,3\},000),\qquad Q(\{2,3\},000).
\]
So \(\SSh(C)=\Sh(C)\) has \(7=|C|\) members: \(C\) is ample and maximum of
VC dimension two.  The three squares are its maximal cubes.  The vertex
\(000\) lies in all three, the vertices \(100,010,001\) in two each, and
\(110,101,011\) in one each.  Hence the corners are \(110,101,011\).

More generally, the maximal cubes of \(Q_n\setminus\{x\}\) are the \(n\)
subcubes of dimension \(n-1\) avoiding \(x\), namely
\(\{y:y_e\ne x_e\}\) for \(e\in[n]\).  A vertex \(y\ne x\) lies in the
\(e\)-th one exactly when \(y_e\ne x_e\), so the corners are the
neighbours of \(x\).  The full cube \(Q_n\) is a single maximal cube, so
all its vertices are corners.
\end{example}

\begin{example}[Why the union condition matters]
\label{ex:non-ample-union}
Let \(\Sigma=\{\varnothing,\{1\},\{2\},\{3\}\}\) on \(E=[3]\), whose facets
are the singletons, and choose the three edges
\[
 Q(\{1\},000)=\{000,100\},\quad Q(\{2\},101)=\{101,111\},\quad
 Q(\{3\},010)=\{010,011\}.
\]
The union has six vertices and shatters the pairs \(\{1,2\},\{1,3\},
\{2,3\}\), which lie outside \(\Sigma\).  Here \(|\Sh(U)|=7>6=|U|\), so the
union is not ample.  Lemma~\ref{lem:facet-cubes} shows that this is the
only way a choice of facet cubes can fail to be ample.
\end{example}

%======================================================================
\section{Facet cubes}
\label{sec:facet}

Call the inclusion-maximal members of \(\Sh(C)\) the \emph{facets} of
\(C\).

\begin{lemma}[Facet cubes]
\label{lem:facet-cubes}
Let \(C\subseteq Q_E\) be a nonempty ample class.
\begin{enumerate}
\item For every facet \(S\) of \(C\), the class \(C\) contains exactly one
      cube with support \(S\), called the \emph{facet cube} of \(S\).
\item The maximal cubes of \(C\) are exactly its facet cubes.
\item \(C\) is the union of its facet cubes, and a vertex of \(C\) is a
      corner if and only if it lies in exactly one facet cube.
\end{enumerate}
Conversely, let \(\Sigma\) be a down-closed family of subsets of \(E\)
and choose one cube with support \(S\) for every inclusion-maximal
\(S\in\Sigma\).  If the union \(U\) of the chosen cubes shatters no set
outside \(\Sigma\), then \(U\) is ample, \(\Sh(U)=\Sigma\), and the chosen
cubes are the facet cubes of \(U\).
\end{lemma}

\begin{proof}
Throughout, \(\Sh(C)=\SSh(C)\) because \(C\) is ample.

(1) Let \(S\) be a facet.  The reduction \(C^S\) is ample by
Theorem~\ref{thm:reductions}.  By \eqref{eq:ssh-reduction}, a set \(U\)
is strongly shattered by \(C^S\) exactly when \(U\cup S\in\SSh(C)\); since
\(S\) is maximal, this holds only for \(U=\varnothing\).  The Sandwich
equality for \(C^S\) gives \(|C^S|=|\SSh(C^S)|=1\).  The members of
\(C^S\) are the placements of the \(S\)-cubes of \(C\), so there is exactly
one.

(2) First, the facet cube \(Q(S,b)\) of a facet \(S\) is maximal: a cube of
\(C\) containing it has a support \(T\supseteq S\) with
\(T\in\SSh(C)=\Sh(C)\), so \(T=S\) and the two cubes coincide.  Conversely,
let \(Q(T,b)\) be a maximal cube of \(C\).  Then \(b\in C^T\).  If \(b\)
had a neighbour \(b'\in C^T\), differing from \(b\) in a coordinate
\(e\notin T\), then \(Q(T,b)\cup Q(T,b')=Q(T\cup\{e\},b\wedge b')\) would
be a larger cube of \(C\).  So \(b\) is an isolated vertex of \(C^T\).
By Theorem~\ref{thm:reductions} the nonempty class \(C^T\) is connected,
hence \(C^T=\{b\}\), and by Theorem~\ref{thm:sandwich}
\(|\SSh(C^T)|=1\), i.e.\ \(\SSh(C^T)=\{\varnothing\}\).  By
\eqref{eq:ssh-reduction} no set \(U\cup T\) with \(U\neq\varnothing\)
lies in \(\SSh(C)=\Sh(C)\).  Hence \(T\) is a facet, and \(Q(T,b)\) is its
facet cube by (1).

(3) Every vertex \(v\in C\) lies in the cube \(Q(\varnothing,v)\), which is
contained in some maximal cube, and by (2) that is a facet cube.  By
definition, \(v\) is a corner exactly when one maximal cube contains it,
which by (2) means exactly one facet cube.

For the converse, every \(T\in\Sigma\) lies in some maximal
\(S\in\Sigma\), and the chosen cube with support \(S\) contains a cube
with support \(T\).  So \(\Sigma\subseteq\SSh(U)\).  By hypothesis
\(\Sh(U)\subseteq\Sigma\).  Theorem~\ref{thm:sandwich} gives
\[
 |\Sigma|\le|\SSh(U)|\le|U|\le|\Sh(U)|\le|\Sigma| ,
\]
so all these numbers are equal.  Hence \(U\) is ample and
\(\SSh(U)=\Sh(U)=\Sigma\).  The facets of \(U\) are the maximal members of
\(\Sigma\), and for each of them the chosen cube is a cube of \(U\) with that
support, hence its facet cube by (1).
\end{proof}

\begin{remark}
Part~(3) does not need the deletion criterion of
Theorem~\ref{thm:corner-deletion}; it follows from the definition of a
corner once the maximal cubes are known.  Part~(2) says in particular that
a maximum class of VC dimension \(d\) has only \(d\)-dimensional maximal
cubes, one for each \(d\)-set of coordinates.  In
Example~\ref{ex:codim-one} the facets are the three pairs, and the
coverage counts \(3,2,2,2,1,1,1\) computed there are the counts of
Lemma~\ref{lem:facet-cubes}(3).
\end{remark}

\begin{corollary}
\label{cor:cornerless-cover}
A nonempty ample class is cornerless if and only if each of its vertices
lies in at least two facet cubes.  In particular a cornerless ample class
has at least two facets, and no facet cube of it is the whole class.
\end{corollary}

There is also a purely local description of corners, valid for every
class.  For \(v\in C\) let \(I_C(v)=\{e\in E: v\oplus e\in C\}\) be the set of
edge directions at \(v\), where \(v\oplus e\) is \(v\) with coordinate \(e\)
flipped.

\begin{lemma}[Local corner criterion]
\label{lem:local-corner}
Let \(C\subseteq Q_E\) and \(v\in C\).  Then \(v\) is a corner of \(C\) if and
only if the cube spanned by \(v\) and its neighbours,
\(Q(I_C(v),\,v\setminus I_C(v))\), is contained in \(C\).
\end{lemma}

\begin{proof}
Every cube of \(C\) through \(v\) has its support inside \(I_C(v)\), since
each coordinate of the support is an edge direction at \(v\).  If the
cube spanned by \(I_C(v)\) at \(v\) lies in \(C\), it therefore contains every
cube of \(C\) through \(v\) and is the unique maximal one.  Conversely, let
\(v\) be a corner with unique maximal cube \(Q\).  For \(e\in I_C(v)\), the
edge \(\{v,v\oplus e\}\) lies in some maximal cube, which is \(Q\); so the
support of \(Q\) contains \(I_C(v)\), and \(Q\) is the spanned cube.
\end{proof}

Hence \(C\) is cornerless if and only if, for every \(v\in C\) and every
\(T\subseteq E\), the implication
\[
 \Bigl(\bigwedge_{e\notin T} v\oplus e\notin C\Bigr)\ \Longrightarrow\
 Q(T,v\setminus T)\not\subseteq C
\]
holds: the case \(T=I_C(v)\) is the lemma, and for \(T\supsetneq I_C(v)\) the
cube \(Q(T,v\setminus T)\) is not contained in \(C\) anyway.  As clauses in the
covering variables this reads
\(\neg c_v\vee\bigvee_{e\notin T}c_{v\oplus e}\vee\bigvee_{w\in
Q(T,v\setminus T)\setminus\{v\}}\neg c_w\).  Replacing clause family~(D)
of Section~\ref{sec:general-formula} by these clauses gives another exact
formula; we compare the two in Section~\ref{sec:open}.

The following two observations reduce the search space.

\begin{lemma}
\label{lem:small-vc}
Let \(C\) be a nonempty cornerless ample class of VC dimension \(d\) on
\(n\) active coordinates.  Then \(3\le d\le n-1\).  If \(C\) is maximum,
then \(3\le d\le n-2\).
\end{lemma}

\begin{proof}
Ample classes of VC dimension at most two have corner peelings
\cite[Proposition~4.11]{ChalopiChepoiMoran22}, and in particular have a
corner, so \(d\ge3\).  If \(d=n\), then \(E\in\Sh(C)\), so \(C=Q_E\) by
ampleness, and every vertex of a cube is a corner.  If \(C\) is maximum and
\(d=n-1\), then \(|C|=2^n-1\) and \(C\) is a cube minus a vertex, which has
corners by Example~\ref{ex:codim-one}.
\end{proof}

We do not use \(d\le n-1\) in the searches for general ample classes; the
formulas for \(d=n\) are included and are trivially refuted.

\subsection{Two reductions}
\label{sec:reductions}

The next two lemmas show that a cornerless ample class with a peripheral
coordinate, or with a product decomposition, yields a cornerless ample class
on fewer coordinates.  A coordinate \(e\) of \(C\) is \emph{peripheral} if,
after possibly exchanging the values \(0\) and \(1\) on \(e\),
\[
 C=(B\times\{0\})\cup(S\times\{1\}),\qquad S\subseteq B,
\]
where \(B=\rho_e^0C\) and \(S=\rho_e^1C\) are viewed in \(Q_{E\setminus\{e\}}\).
Then \(S=C^{\{e\}}\) is the reduction along \(e\).

\begin{lemma}[Peripheral coordinates]
\label{lem:peripheral}
Let \(C\) be ample, let \(e\) be an active peripheral coordinate, and write
\(C=(B\times\{0\})\cup(S\times\{1\})\) with \(S\subseteq B\).  Then \(S\) is a
nonempty ample class, and for every \(s\in S\), the vertex \((s,1)\) is a
corner of \(C\) if and only if \(s\) is a corner of \(S\).  In particular, if
\(C\) is cornerless, then \(S\) is a nonempty cornerless ample class on
\(E\setminus\{e\}\).
\end{lemma}

\begin{proof}
Since \(e\) is active, \(S\neq\varnothing\), and \(S=C^{\{e\}}\) is ample by
Theorem~\ref{thm:reductions}.  Let \(s\in S\).  A cube of \(C\) containing
\((s,1)\) either lies in the layer \(x_e=1\), and then has the form
\(Q\times\{1\}\) with \(Q\) a cube of \(S\) through \(s\), or has \(e\) in its
support, and then has the form \(Q\times\{0,1\}\) with \(Q\subseteq S\cap B=S\).
Conversely, for every cube \(Q\) of \(S\), the cube \(Q\times\{0,1\}\) lies in
\(C\) because \(S\subseteq B\).  Hence the maximal cubes of \(C\) through
\((s,1)\) are exactly the cubes \(Q\times\{0,1\}\) with \(Q\) a maximal cube of
\(S\) through \(s\), and the two corner conditions agree.
\end{proof}

\begin{lemma}[Products]
\label{lem:products}
Let \(C=A\times B\subseteq Q_{E_1}\times Q_{E_2}\) with \(A,B\) nonempty.  The
maximal cubes of \(C\) are the products of maximal cubes of \(A\) and of
\(B\), and \((a,b)\) is a corner of \(C\) if and only if \(a\) is a corner of
\(A\) and \(b\) is a corner of \(B\).  In particular, if \(C\) is cornerless,
then \(A\) or \(B\) is cornerless.
\end{lemma}

\begin{proof}
A cube of \(Q_{E_1\cup E_2}\) with support \(T\) is the product of the cubes
with supports \(T\cap E_1\) and \(T\cap E_2\) through the corresponding
projections, and it lies in \(A\times B\) exactly when both factors lie in
\(A\) and \(B\).  Maximality therefore holds exactly factorwise, and the
maximal cubes through \((a,b)\) are the products of maximal cubes through
\(a\) and through \(b\).  Their number is the product of the two numbers,
which equals \(1\) exactly when both are \(1\).  If \(A\) has a corner \(a\)
and \(B\) has a corner \(b\), then \((a,b)\) is a corner of \(C\).
\end{proof}

Products of ample classes are ample, since shattered and strongly shattered
sets of a product are the unions of those of the factors.  Thus a
cornerless ample product has a cornerless ample factor on fewer
coordinates.

\begin{proposition}
\label{prop:cmin-equals-mamp}
\(\cmin=\mamp\).  More precisely, every ample forbidden pc-minor \(O\) of
\(\Damp\) contains a nonempty cornerless ample subclass, and every nonempty
cornerless ample class \(C\) has an ample forbidden pc-minor of order at
most \(|C|\).
\end{proposition}

\begin{proof}
Let \(O\) be an ample forbidden pc-minor.  Delete corners from \(O\) as
long as one exists.  By Theorem~\ref{thm:corner-deletion} each
intermediate class is ample.  The process cannot empty \(O\), since the
deletion order would then be a corner peeling of \(O\).  It stops at a
nonempty cornerless ample class of order at most \(|O|\).  Hence
\(\cmin\le\mamp\).

Let \(C\) be nonempty, cornerless and ample.  Then \(C\notin\Damp\), since
a corner peeling would begin with a corner.  Among all pc-minors of \(C\)
that are not in \(\Damp\), choose one, \(O\), of least order.  A proper
pc-minor of \(O\) has strictly fewer vertices: a nonempty restriction to a
proper subset deletes vertices, and contracting an active coordinate
identifies the two endpoints of an edge.  Every pc-minor of \(O\) is a
pc-minor of \(C\), so by the choice of \(O\) all proper pc-minors of \(O\)
lie in \(\Damp\).  Since ampleness is preserved by restrictions and
contractions, \(O\) is ample.  Thus \(O\) is an ample forbidden pc-minor,
and \(|O|\le|C|\).  Hence \(\mamp\le\cmin\).
\end{proof}

%======================================================================
\subsection{Unique corners and prescribed endpoints}
\label{sec:unique-corner-search}

A class with a unique corner provides another source of small
obstructions, independent of the contraction-and-lift neighbourhood of a
known example. The construction below is the vertex-gluing rule from
\cite{BousquetDampReader}; we give its short proof to make the size
consequence explicit. We then exclude unique corners in seven coordinates
and deduce Theorem~\ref{thm:E}.

\begin{proposition}
\label{prop:unique-corner-double}
Let $A$ be a nontrivial ample class whose unique corner is $a$.
Identify the corners of two copies of $A$ on disjoint coordinate blocks.
The resulting class is cornerless and ample, with $2|A|-1$ vertices.
Consequently $|A|\ge33$. If such a class exists on seven coordinates,
then $\mamp\le255$.

If $A$ uses $n\ge2$ coordinates, the degree of its unique corner is
at most $n-2$. On seven coordinates this degree is $3$ or $4$.
\end{proposition}
\begin{proof}
Translate $a$ to zero. The shattered supports of the glued class are
precisely those of its two copies: no support using a coordinate from
each block is shattered, since its all-one trace is absent. The two
shadows share just the empty support. Thus their union has
$2|\Sh(A)|-1=2|A|-1$ supports, and the glued class is ample.
Every nonroot vertex has exactly its original maximal cubes, so is
not a corner. At the root an edge from either block exists, since
$A$ is nontrivial and connected. These two edges do not span a square,
so the root is not a corner either. Proposition~\ref{prop:cmin-equals-mamp}
and the bound $\mamp\ge64$ give $2|A|-1\ge64$, hence $|A|\ge33$.
For seven coordinates, $|A|\le128$ gives $\mamp\le255$.

For the degree assertion, a corner of degree $n$ has the full cube as
its corner cube, so cannot be unique. Suppose its degree is $n-1$.
Translate it to zero again, and write $j$ for the only missing incident
direction and $S$ for the other directions. The $j=0$ section is the
full cube $Q_S$. Let $D$ be the $j=1$ section. It does not contain
zero. Each unit vertex in $Q_S$ must have a $j$-neighbour: otherwise
it too is a corner, with unique maximal cube $Q_S$. Hence $D$
contains every unit vertex. The cube complement $Q_S\setminus D$
is ample and connected, by complement closure of ampleness
\cite[Section~3.1]{ChalopiChepoiMoran22}. It contains zero and none
of its neighbours, so it consists of zero alone. Thus $A$ is an
$n$-cube minus one vertex, which has at least two corners for $n\ge2$.
This contradiction proves the assertion.

For the final statement we use the companion results that every ample
family on at most six coordinates peels to any prescribed vertex, and
every nested nonempty ample pair on at most five coordinates peels
relatively \cite{BousquetDampReader}. If the unique corner has degree
one, its incident coordinate has a singleton root section: that
section is ample and connected, and its root is isolated. The other
section uses at most six coordinates and peels to the root's neighbour.
Lift this order, retaining the root, and finish along the last edge.
This contradicts uniqueness of the first corner.

If the degree is two, translate the root to zero and use its two
incident directions as the last coordinates. The sections have the form
\[
 \{(0,00)\}\cup(P\times\{01\})
 \cup(Q\times\{10\})\cup(T\times\{11\}),
\]
where $P,Q,T$ are ample five-coordinate families containing zero.
The first section is a singleton by the same connectedness argument.
Put $R=P\cap T$ and $S=Q\cap T$; these are nonempty ample
reductions of faces. Relatively peel $P$ to $R$ and $Q$ to $S$.
These deletions have no neighbour in either adjacent section, so lift
to the whole class. Next peel $R$ and $S$ to zero, with all mates
retained in $T$. Their nonzero corner cubes lift across the respective
incident edge. Peel $T$ to zero next; its nonzero vertices now have
no mates in the other sections. All these orders exist by the
five-coordinate relative theorem. The remaining square peels to the
original root, again contradicting its being the unique corner.
Together with the degree upper bound, this leaves $3,4,5$.

It remains to exclude degree five. Let $F$ be the five incident
directions. In the other two directions the sections are
$Q_F,P,Q,T$ at $00,10,01,11$, respectively. None of $P,Q,T$
contains zero; for $T$ this follows from isometry and the absence
of the two intermediate neighbours of the root. Every nonzero
vertex of the $00$ cube needs an outside neighbour to avoid being
a corner, so $P\cup Q=Q_F\setminus\{0\}$.
If $P\setminus T$ is nonempty, relative peeling to $P\cap T$
(or ordinary peeling if that intersection is empty) supplies a
corner of $P$ outside $T$. Its cube lifts into the full $00$
section, giving another corner of the whole family. Thus $P\subseteq T$,
and similarly $Q\subseteq T$. Hence $T=P\cup Q=Q_F\setminus\{0\}$.
If one of $P,Q$ is empty, a corner of the $11$ section lifts into
the other section, again contradicting uniqueness. Both are nonempty.

The face $U=(P\times\{0\})\cup(T\times\{1\})$ and the reduction
$K=(P\times\{0\})\cup(Q\times\{1\})$ are ample, with the same
projection $T$. The section-shadow identity gives
$\Sh(T)=\Sh(P)\cup\Sh(Q)$: supports containing the last coordinate
are indexed by $\Sh(P)\cap\Sh(Q)$, so Sandwich equality for $K$
gives equality of cardinalities in the inclusion
$\Sh(P)\cup\Sh(Q)\subseteq\Sh(T)$.
For any ample $Q\subseteq T'\subseteq T$, one has $P\cup T'=T$
and $\Sh(P)\cup\Sh(T')=\Sh(T)$. The family with sections $P,T'$
therefore has $|P|+|T'|$ shattered supports and is ample.
Consequently a five-coordinate relative peeling of $T$ to $Q$
lifts to one of $U$ to $K$. It also lifts into the whole family,
since the deleted words have no mates in the other layer.
The remaining $K$ layer has all its mates in the untouched layer
with sections $Q_F,Q$. Peel $K$ ordinarily using Theorem~\ref{thm:B},
delete the $Q$ layer with all mates in $Q_F$, and finish the cube
at zero. This again contradicts uniqueness of the initial corner.
Only degrees $3$ and $4$ remain.
\end{proof}

The exact facet-cube encoding below also describes this search: require
the root to be covered by exactly one facet cube, and every other
covered vertex by at least two. Fixing the root at zero is harmless,
but translation lex-leaders must then be omitted. Coordinate permutations
still fix the root and can put its incident directions first. This
formulation imposes neither symmetry of the family nor a prescribed
shattered complex.

An alternative exact encoding works directly with membership variables
$c_x$, one for each vertex $x$ of the ambient cube. Lawrence's total
asymmetry characterization \cite[Theorems~3--4]{Lawrence83} says that a
class is ample if and only if every antipodally invariant face section
is empty or full. For each face $F$ of dimension at least two, choose
one representative from each antipodal pair, forming a set $R_F$, and
choose $r_F\in R_F$. Introduce the Boolean variables
$d_{F,x}\leftrightarrow(c_x\mathbin{\mathrm{XOR}}c_{\bar x})$, where
$\bar x$ is the antipode of $x$ in $F$. Impose
\[
 \left(\bigvee_{x\in R_F}d_{F,x}\right)
 \ \lor\ (c_y\leftrightarrow c_{r_F})
 \qquad (y\in R_F\setminus\{r_F\}).
\]
If one antipodal pair differs, these conditions impose nothing further
on that face. If all pairs agree, they force every vertex of the face
to have the same membership value. Thus the formula is exactly total
asymmetry; faces of dimension zero or one need no clauses. A pair of
vertices determines its smallest containing face, so its XOR variable
can be indexed by the pair itself.

For fixed root degree $d$, require the full cube on the first $d$
directions at zero and exclude the other neighbours of zero. To exclude
a corner at any other present vertex $v$, for every direction set $T$
impose the clause
\[
 \neg c_v\ \lor\!!\bigvee_{j\notin T}c_{v^{\{j\}}}
 \ \lor\!\!\bigvee_{\varnothing\ne M\subseteq T}\neg c_{v^M}.
\]
Here $v^M$ denotes flipping the coordinates in $M$. Taking $T$ to be
exactly the incident directions demands that the neighbourhood cube
be incomplete. All other choices either contain a missing neighbour
or omit a present neighbour, so contribute no stronger condition.
This reuses the local corner criterion above and Lawrence's theorem;
it does not assume a prescribed shattered complex. The implementation
\path{scripts/explore_unique_corner_asymmetry.py} checks its ampleness
encoding on every family in $Q_3$ against independent shattered and
strongly shattered counts. Bounded search outcomes are recorded in
the ledger and are not additional nonexistence theorems.

\begin{lemma}[A proper fiber at a root neighbour]
\label{lem:unique-corner-proper-fiber}
Let $A\subseteq Q_D\times Q_R$ be ample with unique corner $(0,0)$,
whose incident directions are exactly $D$, where $D,R$ are nonempty.
Write $A_x=\{y\in Q_R:(x,y)\in A\}$. Every $A_x$ with
$x\ne0$ is a nontrivial ample class containing zero. At least one
unit vertex $e_i\in Q_D$ has $A_{e_i}\ne Q_R$.
\end{lemma}
\begin{proof}
The root cube is $Q_D\times\{0\}$, so every fiber contains zero.
Each fiber is ample, being a face section. If $A_x=\{0\}$ for
$x\ne0$, the incident directions at $(x,0)$ are exactly $D$ and
span the full root cube. This would be another corner.
The fiber $A_0$ is a singleton: it is ample and connected, but zero
has no neighbour in a direction of $R$.

Suppose every unit fiber is full. For each $y\ne0$ in $Q_R$, the
section $B_y=\{x:(x,y)\in A\}$ contains all the unit vertices
of $Q_D$ and excludes zero. Its cube complement is ample, hence
connected, and contains zero but none of its neighbours. It is
therefore the singleton $\{0\}$. Thus
\[
 A=(Q_D\times\{0\})\cup((Q_D\setminus\{0\})\times Q_R).
\]
For any unit $e_i$ and any $y\ne0$, the incident directions at
$(e_i,y)$ are $(D\setminus\{i\})\cup R$. They span the full cube
with coordinate $i$ fixed to one. Hence $(e_i,y)$ is a second
corner, a contradiction.
\end{proof}

\begin{theorem}[The degree-four case is excluded, computer-assisted]
\label{thm:unique-corner-degree-four}
An ample class using seven coordinates cannot have a unique corner of
degree four. Consequently any seven-coordinate ample class with a
unique corner has that corner of degree three.
\end{theorem}
\begin{proof}
Suppose such a class exists, and normalize its root to zero and its
incident directions to the first four coordinates. By
Lemma~\ref{lem:unique-corner-proper-fiber}, some unit root neighbour
has a proper, nontrivial ample fiber in the other three coordinates.
Permute the first four coordinates to make this neighbour $1$, and
permute the last three coordinates to canonicalize its fiber.
No reversal or symmetry of the whole class is imposed.

Encode a fiber $V\subseteq\{0,\ldots,7\}$ by the integer
$m(V)=\sum_{y\in V}2^y$. Checking the odd integers $3\le m<255$
by the shattered and strongly shattered counts gives exactly $62$
proper nontrivial ample fibers containing zero. Under the six
permutations of their three coordinates, their minimum representatives
are
\[
3,7,11,15,23,27,31,43,47,63,127,139,143,171,175,191,239.
\]
Their orbit sizes, in the same order, are
\[
3,3,6,3,1,6,3,3,6,3,1,6,3,3,6,3,3,
\]
which sum to $62$. The routine \texttt{catalog} in
\path{scripts/explore_unique_corner_fibers.py} enumerates every eligible
integer and every coordinate permutation; this is an exhaustive split,
not a sample of fibers.

For each representative, impose its eight membership literals on the
exact total-asymmetry encoding above, together with the unique-root
corner clauses, full coordinate support, and $|A|\ge33$ from
Proposition~\ref{prop:unique-corner-double}. Also require each nonzero
root-cube fiber to be nontrivial, as supplied by
Lemma~\ref{lem:unique-corner-proper-fiber}. Every resulting formula has
$10\,943$ variables and $65\,378$ clauses. All $17$ formulas are
unsatisfiable, with independently checked DRAT proofs. The generating
script is \path{scripts/certify_unique_corner_fibers.py}; formulas,
proofs, checker logs and SHA-256 fingerprints are retained under
\path{evidence/unique_corner_q7_degree4_certificates/}.
Each case was solved afresh with its fiber literals as unit clauses:
the checked proofs do not rely on solver assumptions or clauses learned
in other cases. The split covers every possible proper unit fiber, so
no such class exists. Proposition~\ref{prop:unique-corner-double}
then leaves only degree three.
\end{proof}

\begin{theorem}[The degree-three case is excluded, computer-assisted]
\label{thm:unique-corner-degree-three}
An ample class using seven coordinates cannot have a unique corner of
degree three.
\end{theorem}
\begin{proof}
Normalize the root to zero and its incident directions to the first
three coordinates. Lemma~\ref{lem:unique-corner-proper-fiber} supplies
a proper nontrivial ample four-coordinate fiber at a unit root neighbour.
Permute the root directions to put this neighbour at $1$ and the other
four directions to canonicalize its fiber.

For $V\subseteq\{0,\ldots,15\}$ use the code
$m(V)=\sum_{y\in V}2^y$. Every eligible fiber is tested by enumerating
the odd integers $3\le m<65535$. For each support $T\subseteq[4]$,
its projection is full precisely when the family meets every cell
$\{v:v|_T=b\}$, $b\in\{0,1\}^T$. Count these shattered supports
and retain the code if the count equals its number of vertices.
Equality in the Sandwich bound is exactly ampleness, so this enumerates
all eligible fibers. There are $2763$ labelled fibers. Taking the
minimum code over the $24$ coordinate permutations partitions them into
$185$ orbits. The complete orbit lists are retained in
\path{evidence/unique_corner_q7_degree3_fibers.json}; the routine
\texttt{catalog} in \path{scripts/explore_unique_corner_degree3_fibers.py}
performs this enumeration and independently checks shattered and strongly
shattered counts on each representative.

For each representative, use the same exact membership encoding as in
Theorem~\ref{thm:unique-corner-degree-four}, with a three-dimensional
root cube, nontrivial fibers at the other seven root-cube vertices, and
sixteen unit literals fixing the chosen fiber. The proved lower bound
$|A|\ge33$ is again included. Each formula has $10\,943$ variables and
$65\,379$ clauses. All $185$ fresh formulas have independently checked
DRAT refutations; no solver assumptions or learned clauses from another
case enter the proofs. The complete certificate archive, source snapshots,
logs and fingerprints are in
\path{evidence/unique_corner_q7_degree3_certificates/}, generated by
\path{scripts/certify_unique_corner_degree3_fibers.py}.
The exhaustive orbit cover therefore excludes every possible proper
unit fiber, proving the theorem.
\end{proof}

\begin{proof}[Proof of Theorem~\ref{thm:E}]
The companion prescribed-endpoint theorem covers at most six coordinates
\cite{BousquetDampReader}. In seven coordinates,
Proposition~\ref{prop:unique-corner-double} and
Theorems~\ref{thm:unique-corner-degree-four}
and~\ref{thm:unique-corner-degree-three} exclude a unique corner.
Theorem~\ref{thm:B} supplies at least one corner in every nonempty
ample class. Thus every nontrivial ample class on at most seven
coordinates has at least two corners. Given a prescribed vertex,
delete a different corner. Corner deletion preserves ampleness and
does not increase the number of coordinates. Repeat until only the
prescribed vertex remains. This proves the required peeling order.
\end{proof}

\begin{lemma}[Three fiber cases for a cube of codimension two]
\label{lem:cornerless-codimension-two-fibers}
Let $A\subseteq Q_{D\cup R}$ be cornerless and ample, where $|R|=2$,
and suppose $A$ contains a cube with direction set $D$. After a cube
automorphism, it contains $Q_D\times\{00\}$ and its fiber above
$0\in Q_D$ is one of
\[
 \{00,10\},\qquad \{00,10,01\},\qquad \{00,10,11\}.
\]
Every fiber above $Q_D$ is a nontrivial ample family containing $00$.
\end{lemma}
\begin{proof}
Translate the given cube to $Q_D\times\{00\}$. Each fiber contains
$00$ and is ample by face closure. A singleton fiber would make its
base vertex a corner, since all its incident directions would lie in
the full $D$-cube. Thus every fiber is nontrivial. At least one is
proper, since otherwise $A$ would be the full ambient cube and have
corners. Translate its base index to zero; this preserves the base cube.
A proper nontrivial ample subset of $Q_2$ containing $00$ is either
an edge through $00$, a three-vertex path with $00$ at its centre,
or a three-vertex path with $00$ at an endpoint. The diagonal pair
is not ample. Permuting the two outside coordinates gives the three
listed representatives.
\end{proof}

There is also a relative-peeling obstruction forced by this presentation.
For a hypothetical eight-coordinate example write the four six-coordinate
sections as $N,P,Q,C$, at $00,10,01,11$, where $N=Q_6$. Then
$P,Q$ are proper: otherwise the class contains a seven-cube, whose
one-coordinate fibers must all be nontrivial to avoid corners, forcing
the full eight-cube. Cornerlessness at the base vertices gives $P\cup Q=N$.
If both $P,Q\subseteq C$, then $C=N$, and the cube complement consists
of the two nonempty parts $(N\setminus P)\times\{10\}$ and
$(N\setminus Q)\times\{01\}$, with no edge between them. This contradicts
connectedness of the nonempty ample complement. Thus, after exchanging
the outside coordinates, $P\setminus C\ne\varnothing$.

The intersection $O=P\cap C$ is ample, being the reduction of a face.
It is nonempty: if it were empty, a corner of the nonempty ample class
$P\subseteq Q_6$ would lift with all its mates in the full base section
to a corner of the whole class. The same reasoning shows that no
relative peeling of $P$ to $O$ exists: its first deleted corner would
be outside $C$ and lift to a corner of the whole class. Hence the
hypothetical example forces a strict nonempty ample interval
$O\subsetneq P\subseteq Q_6$ that fails relative peeling, together
with the four-section compatibility above. The prescribed-endpoint
conclusion of Theorem~\ref{thm:E} alone does not exclude such an
interval; relative peeling retains an entire subfamily.

\begin{proposition}[A six-coordinate failure of relative peeling]
\label{prop:relative-six-coordinate-counterexample}
There is a maximum class $P\subseteq Q_6$ of VC dimension three
containing a maximum class $O$ of VC dimension two such that $P$
cannot be corner-peeled to $O$ while retaining every vertex of $O$.
Thus the five-coordinate relative-peeling theorem cannot be extended
to six coordinates, even for maximum classes of consecutive ranks.
\end{proposition}
\begin{proof}
Number the coordinates $0,\ldots,5$ and identify a word with the
integer whose binary digits are its coordinate values. Take
\[
\begin{aligned}
P={}&\{0,1,2,4,5,6,7,8,9,10,11,12,13,14,15,\\
 &16,17,20,21,22,23,25,27,29,31,32,33,40,41,43,\\
 &47,48,49,52,53,54,55,56,57,59,61,63\},\\
O={}&\{0,2,4,6,7,15,16,20,22,23,31,47,48,49,\\
 &52,53,54,55,56,57,61,63\}.
\end{aligned}
\]
Both inclusions $\varnothing\ne O\subsetneq P\subsetneq Q_6$ are
immediate. Direct finite verification gives the following numbers of
shattered supports by cardinality; the strongly shattered counts are
identical:
\[
\begin{array}{c|rrrrrrr|r}
 &0&1&2&3&4&5&6&\text{total}\\\hline
 P&1&6&15&20&0&0&0&42\\
 O&1&6&15&0&0&0&0&22.
\end{array}
\]
Thus $P$ and $O$ are maximum and ample. The complete corner list of
$P$, with the incident coordinate sets, is
\[
\begin{array}{c|c}
\text{corner}&\text{incident directions}\\\hline
2&\{1,2,3\}\\
47&\{2,4,5\}\\
54&\{0,1,5\}\\
56&\{0,3,4\}.
\end{array}
\]
Each of these four vertices belongs to $O$. Consequently no corner
can be deleted while retaining $O$, so the relative peeling cannot start.

The finite checks use only the displayed sets, not the SAT solver
that found them. For each triple, a full cube in $P$ certifies strong
shattering, and for each four-set a missing trace certifies nonshattering.
The corresponding witnesses for $O$ use pairs and triples. The direct
certificate also supplies one missing vertex of the incident cube for
every noncorner of $P$. All witnesses are retained in the accompanying
\texttt{relative\_pair\_q6/direct\_certificate.json} evidence file and
checked by \texttt{verify\_relative\_pair\_q6.py}.
\end{proof}

\begin{theorem}[Rectangular boundaries give forbidden pc-minors]
\label{thm:rectangular-boundary-forbidden}
For $i=1,2$, let $\varnothing\ne K_i\subsetneq A_i$ be ample
classes on disjoint coordinate sets $E_i$, with every coordinate active
in $A_i$. Assume the following:
\begin{enumerate}
 \item $A_i$ and $K_i$ admit complete corner peelings;
 \item every corner $v$ of $A_i$ belongs to $K_i$, but its unique
       maximal cube $M_i(v)$ is not contained in $K_i$;
 \item for each elementary restriction or contraction $\mu$ on $E_i$
       with $\mu K_i\ne\varnothing$, the class $\mu A_i$ peels
       relatively to $\mu K_i$.
\end{enumerate}
Then
\[
 H=(A_1\times K_2)\cup(K_1\times A_2)
\]
is a cornerless ample forbidden pc-minor for $\Damp$. Its order is
\[
 |H|=|A_1||K_2|+|K_1||A_2|-|K_1||K_2|.
\]
\end{theorem}
\begin{proof}
We use the rectangular-boundary shattering formula
\cite[Theorem \texttt{relative-rectangular-boundary-tensor-factorization}]
{BousquetPartitionability26}:
\[
 \operatorname{Sh}(H)=
 (\operatorname{Sh}(A_1)\times\operatorname{Sh}(K_2))
 \cup
 (\operatorname{Sh}(K_1)\times\operatorname{Sh}(A_2)),
\]
where a pair of supports is identified with their disjoint union.
For completeness, if a support in each block is shattered by the larger
class but by neither smaller class, choose one missing trace in each
smaller class. Their combined trace is absent from $H$. In all other
cases the appropriate product slab supplies the witnessing product cube.
Thus the formula is exact and all its supports are strongly shattered,
so $H$ is ample. Inclusion--exclusion gives its order.

Consider a possible corner $(x,y)$ of $H$. If $x\notin K_1$, then
$y\in K_2$ and the section at $y$ is $A_1$. A corner of the whole
class is a corner in this section, so $x$ would be a corner of $A_1$
outside $K_1$, contrary to (2). The symmetric argument handles
$y\notin K_2$. It remains that $x\in K_1$ and $y\in K_2$.
The two coordinate sections are now $A_1,A_2$, so $x,y$ must both
be corners of these classes. All their incident directions occur in
$H$, and the required incident cube is $M_1(x)\times M_2(y)$.
By (2) choose $x'\in M_1(x)\setminus K_1$ and
$y'\in M_2(y)\setminus K_2$. The vertex $(x',y')$ is absent from
$H$, a contradiction. Hence $H$ is cornerless and is not in $\Damp$.

We show directly why a relative peeling in either factor makes the
boundary peelable. Suppose $A_1$ peels to $K_1$. In that relative
order, delete the slice $\{x\}\times K_2$ for each
$x\in A_1\setminus K_1$, using a fixed peeling of $K_2$ inside the
slice. At an intermediate deletion, the incident cube is the product
of the current corner cube at $x$ in the remaining $A_1$ and the
current corner cube in the slice. Later outside slices still contain
all of $K_2$, while slices over $K_1$ contain $A_2\supseteq K_2$,
so this product cube is present. Thus every deletion is valid.
The remainder $K_1\times A_2$ peels by the product of their peelings.

An elementary operation $\mu$ on $E_1$ sends $H$ exactly to
\[
 (\mu A_1\times K_2)\cup(\mu K_1\times A_2).
\]
If $\mu K_1$ is empty, this is the product of two peelable classes,
since peelability is closed under pc-minors. Otherwise (3) and the
preceding slice argument give a peeling; $\mu K_1$ is peelable by
(1) and minor closure. The same argument applies to $E_2$.
Every coordinate is active in $H$, because the smaller classes are
nonempty, so these are proper elementary minors. All proper pc-minors
therefore belong to $\Damp$, proving forbiddenness.
\end{proof}

\begin{corollary}[An explicit rectangular-boundary obstruction]
\label{cor:rectangular-boundary-1364}
Let $O\subset P\subset Q_6$ be the displayed $22/42$ pair in
Proposition~\ref{prop:relative-six-coordinate-counterexample}.
On two disjoint copies of its six coordinates, the class
\[
 H=(P\times O)\cup(O\times P)
\]
is an ample forbidden pc-minor of order $1364$, dimension $12$,
and VC dimension $5$.
\end{corollary}
\begin{proof}
Both factors peel by the prescribed-endpoint theorem through seven
coordinates. Every elementary image of the nested pair uses at most
five coordinates, so the five-coordinate relative-peeling theorem
\cite{BousquetDampReader} supplies condition (3).
The verified corner list of $P$ lies in $O$. Each listed corner has
three incident directions and its full three-cube, whereas $O$ has
VC dimension two, so condition (2) holds. The theorem applies.
Its order is $2\cdot42\cdot22-22^2=1364$. Both coordinate blocks
are active. The shattering formula has maximal support sizes $3+2$
and $2+3$, giving VC dimension five.
\end{proof}

More generally, choose independently, in each six-coordinate block,
a maximum rank-three class $P_i$ and a maximum rank-two subclass
$O_i$ containing all corners of $P_i$. Then
$(P_1\times O_2)\cup(O_1\times P_2)$ is forbidden of the same order,
dimension, and VC dimension. Indeed, every corner of a maximum
rank-three class has a maximal three-cube, so no such cube lies in
$O_i$; the remaining hypotheses follow from the same small-dimension
theorems. Thus this finite input rule needs no separate peeling or
minor tests. The displayed pair proves that its input set is nonempty.

The construction uses no rooted negative input. It converts two
explicit relative failures, with their proper-image peeling conditions,
into an ordinary forbidden pc-minor. The rectangular-boundary operation
and its ampleness are reused from the cited result; the cornerlessness
and minor-minimality criterion above are the added conclusions. This
is a sound generative subfamily, not an exhaustive description of all
forbidden pc-minors. Its $1364$-vertex instance does not improve the
$64$--$267$ minimum-order interval. The independent verifier
\texttt{verify\_rectangular\_obstruction.py} checks its full shattered
family, the absence of every corner, and explicit corner orders for
all $36$ elementary pc-minors.

\begin{proposition}[Intrinsic recovery of a rectangular boundary]
\label{prop:rectangular-boundary-recovery}
Fix a partition of the coordinates of a nonempty class $H$ into two
nonempty blocks $E_1,E_2$. Let $A_1,A_2$ be its projections, and put
\[
 K_1=\{x\in A_1:\{x\}\times A_2\subseteq H\},\qquad
 K_2=\{y\in A_2:A_1\times\{y\}\subseteq H\}.
\]
Suppose $\varnothing\ne K_i\subsetneq A_i$ for $i=1,2$.
Then $H$ has a rectangular-boundary presentation on this split if
and only if
\[
 |H|=|A_1||K_2|+|K_1||A_2|-|K_1||K_2|.
\]
When it exists, this presentation is unique and equals
$(A_1\times K_2)\cup(K_1\times A_2)$.
If $H$ is ample, all four recovered classes are ample. If $H$ is
forbidden for $\Damp$, all four belong to $\Damp$.

Whenever all four recovered classes belong to $\Damp$ and the
presentation exists, $H$ is cornerless if and only if every corner
of each $A_i$ belongs to $K_i$ and its unique maximal cube is not
contained in $K_i$.
\end{proposition}
\begin{proof}
By definition both slabs are contained in $H$. Their union has the
cardinality in the display, so equality of cardinalities is equivalent
to equality of classes. Conversely, in any presentation with nonempty
proper smaller factors, the projections are the larger factors. A
row indexed outside the first smaller factor is exactly the second
smaller factor and hence is not full. Thus the full rows recover the
first smaller factor; full columns recover the second. This proves
uniqueness. The larger factors are projections of $H$. Fixing a word
in $A_2\setminus K_2$ gives the section $K_1$, and symmetrically for
$K_2$. Ampleness follows from minor closure. For a forbidden $H$
these projections and sections are proper pc-minors, so they belong
to $\Damp$.

For necessity of the corner conditions, suppose a corner $x$ of
$A_1$ lies outside $K_1$. Choose a corner $y$ of $K_2$, which exists
because $K_2$ is nonempty and peelable. At $(x,y)$ the incident cube
is the product of the corner cubes in $A_1$ and $K_2$ and lies wholly
in the first slab. This is a corner of $H$, a contradiction. Thus
all corners of $A_1$ lie in $K_1$, and symmetrically for $A_2$.
If a corner cube $M_1(x)$ were wholly in $K_1$, choose any corner
$y$ of $A_2$. It belongs to $K_2$ by what was just proved.
The full incident cube $M_1(x)\times M_2(y)$ lies in the second
slab, again giving a corner. This proves the necessary cube condition
in both blocks. Conversely, the cornerlessness argument in
Theorem~\ref{thm:rectangular-boundary-forbidden} uses exactly these
conditions: outside the two smaller factors a corner cannot project
to an outer corner, and inside them its product cube has a missing
outside--outside vertex. The proper-image hypotheses are not needed
for that argument.
\end{proof}

This is marked recovery and an exact cornerless criterion, not a
converse to the entire forbiddenness theorem: necessity of its relative
peeling conditions on proper images has not been proved. For two
six-coordinate maximum rank-three/rank-two input pairs those conditions
hold automatically, so the stated input rule is exhaustive within
that cornerless rectangular subfamily. There is no assertion that
every cornerless forbidden class has such a split.

For the $1364$-vertex class above, all $2047$ unordered nontrivial
coordinate bipartitions were checked by the recovery criterion.
Exactly one works: the original two six-coordinate blocks, recovering
orders $42,22,42,22$. The script
\texttt{recognize\_rectangular\_boundaries.py} also agrees with direct
enumeration of all possible nonempty proper lower factors on every
nonempty three-coordinate class. The recognition certificate is
retained with the explicit obstruction.

The same construction also supplies a direct small-order target without
maximum assumptions. Any nonempty proper ample pair $K\subset A\subseteq Q_6$
satisfying the corner condition above satisfies the peeling hypotheses
of the construction theorem, by the seven-coordinate endpoint and
five-coordinate relative theorems. Thus its self-pair gives a forbidden
object of order $k(2a-k)$, where $a=|A|$ and $k=|K|$.
Moreover mixing two different admissible pairs cannot improve the minimum
over all self-pairs. If their self-pair orders are
$m_i=k_i(2a_i-k_i)$ and their mixed order is
$m_{12}=a_1k_2+k_1a_2-k_1k_2$, then
\[
 m_{12}^2-m_1m_2=(a_1k_2-a_2k_1)^2\ge0.
\]
All orders are positive, so $m_{12}\ge\min(m_1,m_2)$.
Consequently the entire six-coordinate-input rectangular construction
has an output below $267$ if and only if one admissible pair satisfies
$k(2a-k)\le266$. This is a smaller-obstruction search criterion;
existence of such a pair is not asserted.

\begin{proposition}[Small rectangular inputs have three size profiles]
\label{prop:rectangular-small-input-profiles}
Suppose a nonempty proper ample pair $K\subset A\subseteq Q_6$
retains every corner of $A$, and $k(2a-k)\le266$, where
$a=|A|$ and $k=|K|$. Then $K$ uses all six coordinates and
\[
 (k,a)\in
 \{(7,a):16\le a\le22\}\cup
 \{(8,a):17\le a\le20\}\cup
 \{(9,a):18\le a\le19\}.
\]
In these three cases $K$ has VC dimension at most two and its
shattered family consists of the empty set, all six singletons,
and exactly $k-7$ pairs. The conclusion does not require the
additional corner-cube-escape condition of the construction theorem.
\end{proposition}
\begin{proof}
First, $A$ must use all six coordinates. Otherwise the five-coordinate
relative theorem supplies a first corner outside $K$, a contradiction.
We claim that each side of every active coordinate of $A$ contains
a corner of $A$. Write its two nonempty sections as $A_0,A_1$ and
put $D=A_0\cap A_1$. Connectedness of $A$ implies $D\ne\varnothing$,
and $D$ is ample as a coordinate reduction. If $A_1\setminus D$ is
nonempty, the five-coordinate relative theorem gives a corner of
$A_1$ outside $D$. It has no neighbour across the split coordinate,
so its incident cube is unchanged and it is a corner of $A$.
If $A_1=D$, take any corner of this nonempty five-coordinate class.
Its incident cube doubles across the split coordinate inside $A$,
so its copy in $A_1$ is a corner of $A$. The same argument applies
to $A_0$, proving the claim. Since $K$ retains all corners, it meets
both sides of all six coordinates. Consequently $k\ge7$ by the
Sandwich equality and the six shattered singletons.

The relative peeling theorem for at most eight added concepts
\cite{BousquetDampReader} implies $a-k\ge9$: otherwise this nested
ample pair would admit a first relative deletion. (It is also the
special case of the small graph-core relative theorem where the
whole difference has at most eight vertices.) Therefore
\[
 k(2a-k)\ge k(k+18).
\]
The size bound forces $k\le9$. Substituting $k=7,8,9$ in
$a\ge k+9$ and $a\le\lfloor(266+k^2)/(2k)\rfloor$ gives the stated
ranges. Finally, any shattered triple would add its three pairs
and itself to the seven already shattered empty/singleton supports,
forcing $|K|\ge11$. Thus no triple is shattered, and the Sandwich
equality leaves exactly $k-7$ pair supports.
\end{proof}

\begin{theorem}[Lower bounds for rectangular forbidden classes]
\label{thm:rectangular-input-order-exclusion}
Every class generated by the six-coordinate rectangular-input rule
has at least $272$ vertices. More generally, let a cornerless ample
forbidden pc-minor admit a presentation
\[
 H=(A_1\times K_2)\cup(K_1\times A_2),\qquad
 \varnothing\ne K_i\subsetneq A_i,
\]
on disjoint active coordinate blocks. Then it uses at least twelve
coordinates; if it uses exactly twelve, it has at least $272$ vertices.
The dimension bound is attained by the $1364$-vertex example.
In particular the known twelve-coordinate obstruction of order $267$
has no such rectangular presentation.
\end{theorem}
\begin{proof}[Computer-assisted proof]
We first exclude every strict corner-retaining pair $K\subset A\subseteq Q_6$
whose self-boundary order $k(2a-k)$ is at most $271$.
The proof of Proposition~\ref{prop:rectangular-small-input-profiles}
continues to give exactly the same three profiles when $266$ is
replaced by $271$: $k(k+18)\le271$ forces $7\le k\le9$, and the
upper bounds for $a$ remain $22,20,19$, respectively.

The pre-existing rooted ample inventory on six coordinates was audited
independently through order nine. Starting with $\{0\}$, for every
representative at each order through eight, every possible one-vertex
augmentation was tested by direct shattered-set counting. The resulting
coordinate-permutation orbits coincide exactly with the next inventory
level. Every nonempty ample class on six coordinates peels to any
prescribed vertex, so reversing a peeling ending at zero proves that
this augmentation induction covers every rooted class. Each listed
representative was also checked directly for ampleness. Restricting
to six active coordinates gives $48,162,347$ rooted representatives
at orders $7,8,9$. Taking every translation by a member and every
coordinate permutation removes the root marking and leaves the
following complete cover:
\[
\begin{array}{c|c|c}
 |K|&\text{unrooted lower classes}&\text{allowed }|A|\\\hline
 7&11&16\text{--}22\\
 8&28&17\text{--}20\\
 9&47&18\text{--}19.
\end{array}
\]

For each of these $86$ fixed lower classes, the formula imposes:
$K\subseteq A$; ampleness of $A$ by total asymmetry; the displayed
order bounds; absence of any corner in $A\setminus K$; and an
$A$-neighbour outside $K$ at every corner of $K$. The last condition
is exactly corner-cube escape. Indeed, if a corner of $A$ has its
incident cube in $K$, its neighbours in $A$ and $K$ coincide, making
it a corner of $K$ without an outside neighbour. Conversely such a
corner of $K$ remains a corner of $A$ with the same cube. Thus no
additional restriction is being imposed on the admissible pairs.

All $86$ formulas are unsatisfiable, each with a separately verified
DRAT proof. There are no unresolved cases or unchecked solver answers
in this cover. The exact formulas, proofs, checker logs, source snapshots,
and input inventory are retained under
\path{evidence/rectangular_fixed_lower/} in the ample-corners project.
The program \texttt{prepare\_rectangular\_lower\_cases.py} audits the
inventory; \texttt{certify\_small\_rectangular\_inputs.py} produces and
checks the proofs. A separate coverage audit reproduces all $86$ formula
hashes and checks all $258$ retained artifact hashes. The local escape
equivalence also agrees on every one of the $6050$ nonempty proper
nested pairs on three coordinates. These finite checks supplement the
symbolic coverage and encoding arguments above.

Consequently every admissible self-boundary has order at least $272$.
The mixed-order inequality proved after
Proposition~\ref{prop:rectangular-boundary-recovery} gives the same
lower bound when the two inputs differ.

Finally, for an arbitrary cornerless forbidden $H$ with such a
presentation, Proposition~\ref{prop:rectangular-boundary-recovery}
makes the four input classes peelable and forces the strict
corner-retention conditions. Neither upper class can use at most
five coordinates: relative peeling to its nonempty proper lower
class would supply a corner outside that lower class. Each block
therefore has at least six active coordinates. If their total is
twelve, both have exactly six and the proved order bound applies.
The explicit $1364$ example proves sharpness of the dimension bound.
\end{proof}

This closes the search for a sub-$267$ obstruction within this
rectangular family. It does not raise the global lower bound $64$,
nor does it show that order $272$ is attained in the family. As a
separate direct check, the rectangular recognizer finds no presentation
of the known $267$-vertex class among all $2047$ coordinate splits.
The full generative characterization must therefore include other
structures even at twelve coordinates.

\begin{corollary}[A size criterion for relative peeling on six coordinates]
\label{cor:relative-six-weighted-size}
Let $\varnothing\ne K\subsetneq A\subseteq Q_6$ be ample, with
$k=|K|$ and $a=|A|$. If
\[
 k(2a-k)\le271,
\]
then $A$ peels relatively to $K$. No corner-retention, maximum-class,
or corner-cube-escape hypothesis is required.
\end{corollary}
\begin{proof}
Suppose no such relative peeling exists. Delete corners outside $K$
as long as possible. The process stops at an ample class $B$ properly
containing $K$, with every corner of $B$ in $K$; reaching $K$ would
itself give a relative peeling of $A$.

Whenever a corner $v$ of $B$ has its entire maximal cube contained
in $K$, delete $v$ from both classes. It is also a corner of $K$:
all neighbours of $v$ in $B$ lie in that cube and hence in $K$, so
the incident cubes in the two classes coincide. Both endpoints
therefore remain ample. This is the equal-support case of
common-corner cancellation
\cite[Lemma \texttt{relative-common-corner-cancellation}]
{BousquetPartitionability26}; the relative face set is unchanged.
There is also a direct check that the deletion remains stalled.
Every vertex outside $K$ is nonadjacent to $v$, so its incident
directions do not change on deleting $v$. A missing vertex of its
incident cube cannot become present after a deletion. It therefore
remains a noncorner. The difference of the two classes is unchanged
and nonempty. The smaller class cannot become empty: its last vertex
would have an incident cube consisting only of itself in the larger
class, contradicting connectedness of that larger ample class.

After finitely many such simultaneous deletions, we obtain
$\varnothing\ne K'\subsetneq B'\subseteq Q_6$ such that every
corner of $B'$ lies in $K'$ and every such corner cube escapes $K'$.
The six-coordinate rectangular construction applies to this pair.
Theorem~\ref{thm:rectangular-input-order-exclusion} therefore gives
\[
 |K'|\bigl(2|B'|-|K'|\bigr)\ge272.
\]
But $|B'|\le a$ and $|K'|\le k$, and $y(2x-y)$ is increasing in
each variable on $0<y<x$. Thus the left side is at most
$k(2a-k)\le271$, a contradiction.
\end{proof}

The added conclusion here is a relative-peeling criterion for arbitrary
nested ample pairs of the indicated size. The cancellation operation is
reused from the cited result; it is essential that the two corner cubes
coincide. Deleting an arbitrary common corner with different incident
cubes need not preserve the relative face set. This corollary does not
assert that every pair outside the displayed size range fails to peel.

\begin{lemma}[Nested opposite sections and their cubes]
\label{lem:nested-diagonal-cube-cover}
Let an ample class $X\subseteq Q_n\times Q_2$ have sections
$A,P,Q,C$ at $00,10,01,11$, with $\varnothing\ne C\subseteq A$.
Put $R=C\cap P\cap Q$. Then $R$ is ample, every cube of $C$ is
contained in $P$ or in $Q$, and every connected component of
$C\setminus R$ lies wholly in $P\setminus Q$ or in $Q\setminus P$.
Every corner of $C$ outside $R$ lifts to a corner of $X$.

If $C$ has a corner and $X$ is cornerless, then
$\varnothing\ne R\subsetneq C$, every corner of $C$ lies in $R$,
and $C$ cannot peel relatively to $R$. In particular $X$ has a corner
whenever $C$ has VC dimension at most two or uses at most five active
coordinates.
\end{lemma}
\begin{proof}
Let $K$ be any cube in $C$. Restrict $X$ to the face generated by $K$
and the two last coordinates. Its $00$ and $11$ sections both contain
all of $K$, because $C\subseteq A$. If neither middle section contained
$K$, the nonempty cube complement of this face restriction would be
the union of two nonempty parts at $10$ and $01$, with no edge between
them. This contradicts connectedness of a nonempty ample complement.
Thus $K\subseteq P$ or $K\subseteq Q$. Applying this first to vertices
and then to edges shows $C\subseteq P\cup Q$ and that no edge of
$C\setminus R$ joins its two exclusive sides. This proves the component
assertion. The set $R=A\cap P\cap Q\cap C$ is the double coordinate
reduction of $X$, so is ample.

If $v$ is a corner of $C$ outside $R$, it belongs to exactly one middle
section. Its unique maximal cube $K$ is contained in that section by
the preceding cube-cover argument. Doubling $K$ across the corresponding
last-coordinate edge gives its entire incident cube in $X$, so its
$11$ copy is a corner. If $R=C$, every corner cube of $C$ occurs in
all four sections and extends across both last coordinates, again
giving a corner of $X$.

Consequently, if $X$ is cornerless and $C$ has a corner, then $R$ is
nonempty and proper and contains every corner of $C$. No first relative
deletion is possible. Finally, relative peeling holds for every nested
ample pair with upper VC dimension at most two, and for every pair
using at most five coordinates \cite{BousquetDampReader}; these classes
also have corners. These results exclude the required stalled interval
in either stated case.
\end{proof}

\begin{proposition}[Exact compatibility over a full base]
\label{prop:full-base-shadow-defect}
Let $P,Q,C\subseteq Q_n$ be arbitrary, including possibly empty
families, and let $X\subseteq Q_n\times Q_2$ have sections
$Q_n,P,Q,C$ at $00,10,01,11$. For any family $H$, put
$\delta(H)=|\operatorname{Sh}(H)|-|H|$, and define
\[
\begin{split}
 E_P&=\operatorname{Sh}(P\cup C)\setminus
       (\operatorname{Sh}(P)\cup\operatorname{Sh}(C)),\\
 E_Q&=\operatorname{Sh}(Q\cup C)\setminus
       (\operatorname{Sh}(Q)\cup\operatorname{Sh}(C)),\\
 E_C&=\operatorname{Sh}(C)\setminus
       (\operatorname{Sh}(P)\cup\operatorname{Sh}(Q)).
\end{split}
\]
Then the following identity is a sum of nonnegative terms:
\begin{equation}
 \delta(X)=\delta(P)+\delta(Q)+\delta(C)
                  +|E_P|+|E_Q|+|E_C|.
 \label{eq:full-base-shadow-defect}
\end{equation}
Consequently $X$ is ample if and only if $P,Q,C$ are ample and
$E_P=E_Q=E_C=\varnothing$.

Suppose these conditions hold. For a corner $v$ of an ample family
$H$, write $K_H(v)$ for its unique maximal cube, viewed as a set of
words in $Q_n$, and put $R=P\cap Q\cap C$. Then $X$ is cornerless
if and only if all the following conditions hold:
\begin{enumerate}
 \item $P\cup Q=Q_n$, with $P\ne Q_n$ and $Q\ne Q_n$;
 \item every corner $v$ of $P$ lies in $C$ and satisfies
       $K_P(v)\not\subseteq Q\cap C$;
 \item every corner $v$ of $Q$ lies in $C$ and satisfies
       $K_Q(v)\not\subseteq P\cap C$;
 \item every corner $v$ of $C$ lies in $R$ and satisfies
       $K_C(v)\not\subseteq R$.
\end{enumerate}
\end{proposition}
\begin{proof}
Write $\mathcal A=\operatorname{Sh}(P)$,
$\mathcal B=\operatorname{Sh}(Q)$, and
$\mathcal D=\operatorname{Sh}(C)$. Splitting shattered supports
according to their intersection with the two last coordinates gives
\[
 |\operatorname{Sh}(X)|=2^n+
 |\operatorname{Sh}(P\cup C)|+
 |\operatorname{Sh}(Q\cup C)|+
 |\mathcal A\cap\mathcal B\cap\mathcal D|.
\]
Indeed, a support using just the first of the two last coordinates
requires every old trace in $P\cup C$; its other section is full.
The second coordinate is symmetric. A support using both requires
every old trace separately in $P,Q,C$.
For any three finite sets one has
\[
 |\mathcal A\cup\mathcal D|+|\mathcal B\cup\mathcal D|
 +|\mathcal A\cap\mathcal B\cap\mathcal D|
 =|\mathcal A|+|\mathcal B|+|\mathcal D|
   +|\mathcal D\setminus(\mathcal A\cup\mathcal B)|.
\]
Substitute this equality and $|X|=2^n+|P|+|Q|+|C|$ to obtain
\eqref{eq:full-base-shadow-defect}. The Sandwich inequality makes
each $\delta$ nonnegative, with equality exactly for ample families.
This proves the ampleness criterion, including empty sections.

For corners, inspect the cube spanned by all incident directions.
At a base vertex, the old directions span the full $Q_n$. If neither
middle neighbour is present, this vertex is a corner. If exactly one
is present, it is a corner precisely when that middle section is
$Q_n$. If both are present, it is a corner precisely when all three
other sections are full. Moreover a cornerless class cannot have
$P=Q_n$ or $Q=Q_n$: it would contain a full codimension-one cube,
and avoiding a corner at each vertex of that cube forces every
opposite vertex to be present, giving the full cube and a contradiction.
These observations give exactly condition (1).

At $(v,10)$ the base neighbour is always present. A corner there
requires $v$ to be a corner of $P$. If $v\notin C$, its incident cube
is automatically complete; if $v\in C$, that cube is complete exactly
when $K_P(v)\subseteq Q\cap C$. This gives (2), and interchanging
the last coordinates gives (3).

At $(v,11)$ a corner requires $v$ to be a corner of $C$. Lemma~\ref{lem:nested-diagonal-cube-cover} shows that a corner outside $R$ lifts
to a corner of $X$. For a corner in $R$, both middle neighbours are
present, and its incident cube is complete exactly when
$K_C(v)\subseteq P\cap Q$; the base is full. Since $K_C(v)\subseteq C$,
this is equivalent to containment in $R$, proving (4).
Every vertex lies in one of these four sections, so the conditions
are also sufficient.
\end{proof}

This is an exact reduction of the full-base search to three families
on two fewer coordinates, with no maximum-class assumption. It
specializes the general counting of shattered supports by section
indices: the full base makes every remaining defect nonnegative
and gives the explicit compatibility conditions above. It does not
classify the triples satisfying those conditions. In particular it
is not an exhaustive generator for all forbidden pc-minors. The
identity and the local corner formulas are independently checked
on all $4096$ labelled triples of subsets of $Q_2$, including the
$819$ ample full-base families, by
\texttt{verify\_full\_base\_defect.py}.

\begin{corollary}[A single edge-fiber normalization]
\label{cor:full-base-edge-normalization}
Let a cornerless ample class have sections $Q_n,P,Q,C$.
Then both $P\setminus(Q\cup C)$ and $Q\setminus(P\cup C)$ are
nonempty. In particular, translating only the old coordinates puts
its fiber over zero in the form $\{00,10\}$. This one fiber type
suffices to cover all cornerless full-base classes up to translation.
\end{corollary}
\begin{proof}
The exact corner criterion makes $P,Q$ proper and $P\cup Q=Q_n$.
Also $C\ne Q_n$: otherwise Lemma~\ref{lem:nested-diagonal-cube-cover},
applied to the full cube $C$, would put $Q_n$ in $P$ or $Q$.
The ampleness criterion gives
\[
 \operatorname{Sh}(Q\cup C)
 =\operatorname{Sh}(Q)\cup\operatorname{Sh}(C).
\]
If $Q\cup C=Q_n$, the full coordinate support would be shattered
by $Q$ or $C$, forcing that family to be $Q_n$, a contradiction.
Hence some $x$ lies outside $Q\cup C$. Since $P\cup Q=Q_n$, it lies
in $P$, so its four-section fiber is exactly $\{00,10\}$.
The symmetric argument proves $Q\setminus(P\cup C)\ne\varnothing$.
Translation by $x$ preserves the full base and gives the normalization.
\end{proof}

Thus the earlier three-case fiber cover can be replaced by a single
case. This does not assert that the other fiber types cannot occur at
other old vertices. It chooses one guaranteed edge fiber, without
assuming invariance of the family under translations or permutations.

\begin{proposition}[A certified degree restriction at an exclusive edge fiber]
\label{prop:full-base-exclusive-degree}
Let $X\subseteq Q_6\times Q_2$ be cornerless and ample, with sections
$Q_6,P,Q,C$ at $00,10,01,11$. Every vertex of
$P\setminus(Q\cup C)$ has at least four neighbors in $P$.
Symmetrically, every vertex of $Q\setminus(P\cup C)$ has at least
four neighbors in $Q$.
\end{proposition}
\begin{proof}
Choose $x\in P\setminus(Q\cup C)$ and translate the six old
coordinates so that $x=0$. The exact corner criterion implies that
every corner of $P$ belongs to $C$. Thus $0$ is not a corner of $P$;
in particular its degree in $P$ is at least two, since a vertex of
degree zero or one belongs to a unique maximal cube.
Permuting the six old coordinates preserves the full base and the
fiber over zero. If the degree is $k$, we may therefore require
that its neighbors are exactly the first $k$ unit vectors.

For $k=2,3$, the exact full-base shadow encoding with
$0\in P\setminus(Q\cup C)$ and these six membership literals is
unsatisfiable. Each formula has $9629$ variables and $81747$ clauses.
The formulas were reconstructed byte-for-byte, and their DRAT
refutations independently checked locally after the remote checks.
The formulas, compressed proofs, hashes and replay records are in
\texttt{papers/ample-corners/evidence/q8\_full\_base\_remote};
\texttt{verify\_q8\_full\_base\_remote.py} reproduces the checks.
This rules out degrees two and three. Exchanging the two new
coordinates proves the symmetric statement.
\end{proof}

The degree-four, degree-five and degree-six cases remain unresolved:
the remote solver reached its $600$-second limit in each case, as did
two unsplit encodings. The proposition applies only to the indicated
exclusive fibers. It neither excludes all six-cube classes nor changes
the interval $[64,267]$ for the smallest ample obstruction.

\begin{proposition}[Triple shattering and a small overlap force a corner]
\label{prop:maximum-relative-pair-no-completion}
Suppose an ample class $X\subseteq Q_n\times Q_2$ has sections
$Q_n,P,Q,C$ at $00,10,01,11$, with $C\ne\varnothing$.
If $P$ shatters every three-coordinate support and
$\operatorname{VCdim}(P\cap C)\le2$, then $X$ has a corner.
In particular this excludes every maximum rank-three side with a
maximum rank-two overlap, without requiring either maximum hypothesis
in the general statement.
\end{proposition}
\begin{proof}
Ampleness of the face with sections $P,C$ gives
$\operatorname{Sh}(P)\cap\operatorname{Sh}(C)
=\operatorname{Sh}(P\cap C)$: a support shattered in both sections,
with the separating coordinate adjoined, is strongly shattered in the
face. Since every triple belongs to $\operatorname{Sh}(P)$ and no
triple belongs to the right side, $C$ has VC dimension at most two.
Lemma~\ref{lem:nested-diagonal-cube-cover} now supplies a corner of $X$.
\end{proof}

\begin{corollary}[A connected difference excludes a maximum opposite section]
\label{cor:connected-maximum-opposite-exclusion}
Let a cornerless ample class have sections $Q_n,P,Q,C$, and put
$R=C\cap P\cap Q$. Suppose $C$ is maximum of VC dimension $d\ge1$
and $R$ is maximum of VC dimension $d-1$. Then the induced graph
on $C\setminus R$ is disconnected.
\end{corollary}
\begin{proof}
If this graph were connected, the component assertion of
Lemma~\ref{lem:nested-diagonal-cube-cover} would place it entirely
in one exclusive side. After interchanging the middle coordinates,
$C\subseteq P$ and $C\cap Q=R$. The section-shadow identity for the
face with sections $Q,C$ gives
$\operatorname{Sh}(Q)\cap\operatorname{Sh}(C)=\operatorname{Sh}(R)$.
All $d$-sets lie in $\operatorname{Sh}(C)$ and none in
$\operatorname{Sh}(R)$, so $\operatorname{VCdim}(Q)\le d-1$.
Since $Q$ contains the maximum class $R$, the Sauer bound forces $Q=R$.

Cornerlessness at the full-base vertices forces $P\cup Q=Q_n$:
a base vertex with neither outside neighbour is a corner. As
$Q=R\subseteq C\subseteq P$, this gives $P=Q_n$.
The class then contains a full $(n+1)$-cube. Every one-coordinate
fiber above this cube must be nontrivial to avoid a corner at its
base, forcing the full $(n+2)$-cube, which has corners. This is the
required contradiction.
\end{proof}

The $42/22$ pair above is excluded in both marked positions. The
proposition excludes placing its larger class in a middle section
with the smaller class as its overlap with $C$. The corollary also
excludes placing the larger class at $11$ and the smaller class at
$R$: their $20$-vertex difference is connected. A spanning tree is
included in the direct finite certificate for the pair. These statements
concern the indicated full-base completions, not every use of the pair
in a larger construction.

Proposition~\ref{prop:maximum-relative-pair-no-completion} explains
why the pair in Proposition~\ref{prop:relative-six-coordinate-counterexample}
does not realize the four-section obstruction. Fixing its $P$
and imposing $P\cap C=O$, with $P\cup Q=Q_6$, leaves the other
memberships of $Q,C$ free. The exact cornerless-ampleness formula is
unsatisfiable, with an independently checked DRAT proof retained in
\path{evidence/relative_pair_q6_completion/}. This excludes the marked
completion of this particular pair, not the other six-coordinate
relative intervals or the eight-coordinate six-cube case.

For the eight-coordinate cornerless problem, VC dimension six implies
the presence of a six-cube, so this lemma gives an exhaustive three-case
split. The script \path{scripts/explore_cornerless_q8_large_cube.py}
combines it with the exact total-asymmetry and local corner formulas.
These are necessary normalization choices, not assumptions that the
whole family is invariant under a symmetry. All three initial bounded
runs were inconclusive; their formulas and records are retained under
\path{evidence/cornerless_q8_six_cube/}. In particular this split does
not exclude the eight-coordinate VC-dimension-six case.

%======================================================================
\begin{proposition}[Supports required opposite a degree-two corner]
\label{prop:square-opposite-support-bound}
Let $A,B,C\subseteq\{0,1\}^n$ contain zero, where $n\le7$, and suppose
\[
 Y=\{(0,00)\}\cup(A\times\{01\})\cup(B\times\{10\})
       \cup(C\times\{11\})
\]
is ample. Set $K=A\cap C$ and define
\[
 \mathcal G=\operatorname{Sh}(A)\setminus\operatorname{Sh}(K),
 \qquad
 \mathcal L=\{T\subseteq[n]:\text{ no }S\in\mathcal G
                                      \text{ satisfies }S\subseteq T\}.
\]
Then
\[
 |B\setminus C|\le |\mathcal L\setminus\operatorname{Sh}(A)|.
\]
If $(0,00)$ is the unique corner of $Y$, necessarily
\[
 |\mathcal L\setminus\operatorname{Sh}(A)|\ge9.
\]
The same conclusion holds after exchanging $A$ and $B$.
\end{proposition}
\begin{proof}
The coordinate faces with sections $(A,C)$ and $(B,C)$ are ample.
Their section--reduction identities give
\[
 \operatorname{Sh}(A\cap C)=\operatorname{Sh}(A)\cap\operatorname{Sh}(C),
 \qquad
 \operatorname{Sh}(B\cap C)=\operatorname{Sh}(B)\cap\operatorname{Sh}(C).
\]
For completeness, the first identity follows from shadow containment
and equality of cardinalities: the ample face and its projection imply
$\operatorname{Sh}(A\cup C)=\operatorname{Sh}(A)\cup\operatorname{Sh}(C)$;
then $|A\cap C|=|A|+|C|-|A\cup C|$ is the size of the shadow
intersection. The reduction $A\cap C$ is ample. The second identity
is proved in the same way.
By the unrestricted degree-two-section shadow criterion of
\cite{BousquetDampReader},
$\operatorname{Sh}(A)\cap\operatorname{Sh}(B)
 \subseteq\operatorname{Sh}(C)$.
Thus $\operatorname{Sh}(B)$ contains no member of $\mathcal G$.
As a down-set it contains no superset of such a member either, so
$\operatorname{Sh}(B)\subseteq\mathcal L$. Moreover,
\[
 \operatorname{Sh}(B)\setminus\operatorname{Sh}(C)
       \subseteq\mathcal L\setminus\operatorname{Sh}(A).
\]
The second section--reduction identity and ampleness now give
\[
 |B\setminus C|
 =|\operatorname{Sh}(B)\setminus\operatorname{Sh}(B\cap C)|
 =|\operatorname{Sh}(B)\setminus\operatorname{Sh}(C)|,
\]
proving the size bound.

Suppose the right-hand side of the bound is at most eight.
If $B\setminus C$ is nonempty, the relative peeling theorem for at
most eight added vertices, proved in the argument of
the small graph-core relative-peeling corollary of
\cite{BousquetDampReader}, supplies a corner
of $B$ outside $B\cap C$. This corner is nonzero and has no neighbour
in either adjacent section $C$ or $\{0\}$. Its copy in $B\times\{10\}$
is therefore a corner of $Y$ distinct from $(0,00)$.
If $B\subseteq C$ and $|B|>1$, the prescribed-endpoint theorem
(Theorem~\ref{thm:E}) supplies a peeling
of $B$ ending at zero, whose first vertex is a nonzero corner.
Its incident cube doubles into $C$, again giving a corner of $Y$.
Finally, if $B=\{0\}$, the vertex $(0,10)$ is a corner: its only
incident directions are the two new coordinates and their full square
is present. These cases exhaust the possibilities and prove the claim.
\end{proof}

\begin{example}[A stalled pair excluded by the support bound]
\label{ex:square-thirty-eight-pair}
Encode a vertex of $\{0,1\}^6$ by the integer whose binary expansion
lists its coordinates. Let
\begin{align*}
 K={}&\{0,1,3,4,5,8,9,10,11,12,13,20,21,28,29,40,42,44,46,52,53,58,60,62\},\\
 A={}&K\cup\{16,17,32,33,34,35,36,38,41,43,48,49,50,54\}.
\end{align*}
Direct trace and cube enumeration verifies that these classes are ample,
with shattered-support counts $(1,6,15,16)$ and $(1,6,15,2)$ for $A$
and $K$, respectively. The corners of $A$ are $3,10,29,53,58$, all
in $K$, so relative peeling cannot start. Nevertheless this pair cannot
occur as $(A,A\cap C)$ in a square with a unique degree-two corner.
Indeed, $\mathcal G$ consists of fourteen triples and meets the set of
triples of every four-element support. Hence $\mathcal L$ has no support
of size four or more, and
\[
 \mathcal L\setminus\operatorname{Sh}(A)=\{7,19,38,50\}
\]
in the same binary encoding of supports. The preceding proposition
applies with bound four. The finite checks, including a member of
$\mathcal G$ in every four-element support, are reproduced by
\path{scripts/verify_square_support_bound.py} in the ample-corners project.
This rules out this marked completion, not every use of the pair in
an obstruction construction.
\end{example}

\section{Exact propositional encodings}
\label{sec:encoding}

Lemma~\ref{lem:facet-cubes} replaces the global notion of ampleness by
three local requirements: a down-closed complex, one cube per facet, and a
family of forbidden traces.  Each is a set of clauses.

\subsection{The general formula}
\label{sec:general-formula}

Fix \(n\ge1\) and \(1\le d\le n\), and let \(E=[n]\).  For subsets we use
the words of \(Q_E\) as above.  The formula \(F_{n,d}\) has the following
variables.
\begin{itemize}
\item \(\sigma_T\) for \(T\subseteq E\): \(T\) belongs to the complex
      \(\Sigma\);
\item \(\varphi_T\) for \(T\subseteq E\): \(T\) is a facet of \(\Sigma\);
\item \(s_{T,b}\) for \(T\subseteq E\) and \(b\subseteq E\setminus T\):
      the cube \(Q(T,b)\) is chosen;
\item \(c_v\) for \(v\in Q_E\): \(v\) is covered by a chosen cube;
\item \(y_{T,p}\) for \(\varnothing\neq T\subseteq E\) and \(p\subseteq
      T\): the trace \(p\) is missing on \(T\);
\item auxiliary variables of the cardinality encodings below.
\end{itemize}
Its clauses are:
\begin{enumerate}
\item[(\(\Sigma\))] \(\sigma_\varnothing\); \(\sigma_{\{i\}}\) for every \(i\in E\);
      and \(\neg\sigma_T\vee\sigma_{T\setminus\{i\}}\) for all \(i\in T\).
\item[(VC)] \(\bigvee_{|T|=d}\sigma_T\), and \(\neg\sigma_T\) for every
      \(T\) with \(|T|=d+1\).
\item[(\(\Phi\))] \(\varphi_T\leftrightarrow\bigl(\sigma_T\wedge
      \bigwedge_{i\notin T}\neg\sigma_{T\cup\{i\}}\bigr)\).
\item[(S)] \(\varphi_T\leftrightarrow\bigvee_b s_{T,b}\), and at most one
      \(s_{T,b}\) is true for each \(T\).
\item[(C)] \(c_v\leftrightarrow\bigvee_{T\subseteq E}s_{T,\,v\setminus T}\).
\item[(D)] \(c_v\rightarrow\) at least two of the literals
      \(s_{T,\,v\setminus T}\), \(T\subseteq E\), are true.
\item[(Y)] For every \(T\neq\varnothing\):
      \(\sigma_T\vee\bigvee_{p\subseteq T}y_{T,p}\), and
      \(\neg y_{T,\,v\cap T}\vee\neg c_v\) for every \(v\in Q_E\).
\end{enumerate}
The constraints ``at most one'' in~(S) and ``at least two'' in~(D) are
written with sequential counters \cite{Sinz05}; any correct encoding would
do.  Clause family~(C) uses the fact from Section~\ref{sec:prelim} that
\(v\in Q(T,b)\) if and only if \(b=v\setminus T\).  Clause family~(Y) says:
if \(T\notin\Sigma\), then some trace \(p\) on \(T\) is realized by no
covered vertex.

\begin{theorem}[Exactness]
\label{thm:encoding}
The formula \(F_{n,d}\) is satisfiable if and only if there is a nonempty
cornerless ample class \(C\subseteq Q_{[n]}\) of VC dimension exactly
\(d\) in which every coordinate is active.  More precisely, for every
satisfying assignment the set \(U=\{v:c_v\}\) is such a class, with
\(\Sh(U)=\{T:\sigma_T\}\) and with \(\{Q(T,b):s_{T,b}\}\) as its facet
cubes; and every such class arises in this way.
\end{theorem}

\begin{proof}
Let an assignment satisfy \(F_{n,d}\), and let
\(\Sigma=\{T:\sigma_T\}\).  By~(\(\Sigma\)), \(\Sigma\) is down-closed and contains
all singletons.  By~(\(\Phi\)), \(\varphi_T\) holds exactly for the maximal
members of \(\Sigma\).  By~(S), exactly one cube \(Q(T,b_T)\) is chosen
for each maximal \(T\in\Sigma\) and none for any other \(T\).  By~(C), \(U\)
is the union of the chosen cubes.  By~(Y), for each \(T\notin\Sigma\),
some trace \(p\) on \(T\) satisfies \(y_{T,p}\), and then no \(v\in U\)
has \(v\cap T=p\); so \(U\) does not shatter \(T\).  The converse half of
Lemma~\ref{lem:facet-cubes} shows that \(U\) is ample with
\(\Sh(U)=\Sigma\) and that the chosen cubes are its facet cubes.  It is
nonempty and all coordinates are active because \(\Sigma\) contains the
singletons.  By~(VC) and down-closedness, its VC dimension is exactly
\(d\).  By~(D), every vertex of \(U\) lies in at least two chosen cubes,
that is, in two facet cubes, so \(U\) is cornerless by
Corollary~\ref{cor:cornerless-cover}.

Conversely, let \(C\) be such a class.  Set \(\sigma_T\) true iff
\(T\in\Sh(C)\), \(\varphi_T\) true iff \(T\) is a facet, \(s_{T,b}\) true iff
\(Q(T,b)\) is a facet cube, and \(c_v\) true iff \(v\in C\).  For
\(T\notin\Sh(C)\), pick one trace \(p\) on \(T\) that \(C\) does not
realize and set \(y_{T,p}\) true; set all other \(y\) false.  Clauses~(\(\Sigma\)),
(VC), (\(\Phi\)) and~(Y) hold by construction.  Clause~(S) holds by
Lemma~\ref{lem:facet-cubes}(1).  Clause~(C) holds by
Lemma~\ref{lem:facet-cubes}(3), and~(D) by
Corollary~\ref{cor:cornerless-cover}.  The auxiliary variables of the
cardinality encodings can be set accordingly.
\end{proof}

\begin{remark}
Deleting~(D) gives a formula whose models are exactly the ample classes of
VC dimension \(d\) with all coordinates active.  We used this variant as a
sanity check: on small instances it is satisfiable, and every model
returned was confirmed ample by the direct test \(|\Sh(U)|=|\SSh(U)|=|U|\).
\end{remark}

\subsection{The maximum-class formula}
\label{sec:maximum-formula}

For maximum classes the complex is known in advance: \(\Sigma\) consists of
all sets of size at most \(d\).  The formula \(M_{n,d}\) has variables
\(s_{T,b}\) only for \(|T|=d\), covering variables \(c_v\), and missing-trace
variables \(y_{T,p}\) only for \(|T|=d+1\).  Its clauses are: exactly one
\(s_{T,b}\) is true for each \(d\)-set \(T\); clauses~(C) and~(D) with \(T\)
ranging over \(d\)-sets; and clauses~(Y) for every \((d+1)\)-set \(T\),
with \(\sigma_T\) omitted.

\begin{proposition}
\label{prop:maximum-encoding}
The formula \(M_{n,d}\) is satisfiable if and only if there is a
cornerless maximum class of VC dimension \(d\) on \([n]\).
\end{proposition}

\begin{proof}
Given a model, let \(U\) be the union of the chosen \(d\)-cubes.  It
strongly shatters every set of size at most \(d\), and by~(Y) it shatters
no \((d+1)\)-set, hence no larger set.  With \(\Sigma\) the sets of size at
most \(d\), the converse half of Lemma~\ref{lem:facet-cubes} shows that
\(U\) is ample with \(\Sh(U)=\Sigma\), hence maximum, and that the chosen
cubes are its facet cubes.  By~(D) and
Corollary~\ref{cor:cornerless-cover}, \(U\) is cornerless.  Conversely, a
cornerless maximum class \(C\) is ample with \(\Sh(C)\) equal to the sets
of size at most \(d\) (Section~\ref{sec:prelim}); its facet cubes, its
vertex set, and one missing trace per \((d+1)\)-set give a model, as in
Theorem~\ref{thm:encoding}.
\end{proof}

\subsection{Lifts of a fixed base}
\label{sec:lift-formula}

Let \(B\subseteq Q_E\) be fixed, let \(z\notin E\) be a new coordinate,
and let \(N\) be an integer.  A \emph{lift} of \(B\) is a class
\(Y\subseteq Q_{E\cup\{z\}}\) whose contraction \(\pi_zY\) is \(B\), that is,
\(Y=(P\times\{0\})\cup(Q\times\{1\})\) with \(P\cup Q=B\).  Its vertices lie
in the candidate set \(V=B\times\{0,1\}\).

The formula \(L_{B,N}\) is obtained from \(F_{|E|+1,\cdot}\) on the
coordinate set \(E\cup\{z\}\) with three changes: clause~(VC) is omitted;
cube variables \(s_{T,b}\) are created only for cubes \(Q(T,b)\subseteq V\),
and covering variables \(c_v\) only for \(v\in V\), with~(C), (D) and~(Y)
restricted accordingly (a vertex outside \(V\) is never covered); and two
constraints are added,
\[
 c_{(x,0)}\vee c_{(x,1)}\quad(x\in B),\qquad
 \sum_{v\in V}c_v\le N,
\]
the second written with a totalizer \cite{BailleuxBoufkhad03}.  If a facet
of \(\Sigma\) has no cube inside \(V\), clause~(S) forces \(\varphi_T\) to
be false.

\begin{proposition}
\label{prop:lift-encoding}
The formula \(L_{B,N}\) is satisfiable if and only if \(B\) has a
cornerless ample lift \(Y\) with \(|Y|\le N\) in which every coordinate of
\(E\cup\{z\}\) is active.
\end{proposition}

\begin{proof}
The proof of Theorem~\ref{thm:encoding} applies verbatim.  The restriction
to cubes inside \(V\) loses nothing: a facet cube of a lift \(Y\) lies in
\(Y\subseteq V\).  The added clauses say exactly that
\(\pi_zY\supseteq B\), hence \(\pi_zY=B\) since \(Y\subseteq V\), and that
\(|Y|\le N\).
\end{proof}

\subsection{Symmetry breaking}
\label{sec:lex}

The hyperoctahedral group \(G_n\) of cube automorphisms, generated by
coordinate permutations and coordinate reflections \(x\mapsto x\oplus e_j\),
acts on classes and preserves ampleness, the VC dimension, active
coordinates, facets, facet cubes and corners.  It acts on the cube
variables by \(g\cdot s_{T,b}=s_{g(Q(T,b))}\), where \(g(Q(T,b))\) is again
a cube, with support \(g(T)\).

Fix a total order on the cube variables.  For an assignment \(s\) of the
cube variables and \(g\in G_n\), let \(s^g\) be the assignment
\(s^g_{T,b}=s_{g(Q(T,b))}\).  For a set \(R\) of generators, the
\emph{lex-leader constraints} require \(s\le_{\mathrm{lex}}s^g\) for every
\(g\in R\) \cite{CrawfordGinsbergLuksRoy96}.  They are encoded by the
standard chain of equality variables: at the first position where \(s\)
and \(s^g\) differ, \(s\) has \(0\) and \(s^g\) has \(1\).  We use for \(R\)
the \(n-1\) adjacent transpositions and the \(n\) single-coordinate
reflections, all of which are involutions.

\begin{proposition}
\label{prop:lex-sound}
Adding the lex-leader constraints for \(R\) to \(F_{n,d}\) or to \(M_{n,d}\)
preserves satisfiability.
\end{proposition}

\begin{proof}
By Theorem~\ref{thm:encoding} (resp.\ Proposition~\ref{prop:maximum-encoding}),
models correspond to classes \(C\), and the cube variables of the model are
the indicator vector \(s(C)\) of the facet cubes of \(C\), which determines
\(C\) as their union.  For \(g\in G_n\) the class \(g^{-1}(C)\) is again
admissible and \(s(g^{-1}(C))=s(C)^g\), because \(g\) maps facet cubes of
\(g^{-1}(C)\) to facet cubes of \(C\).  Choose \(C\) in its \(G_n\)-orbit
with \(s(C)\) lexicographically least.  Then \(s(C)\le_{\mathrm{lex}}
s(C)^g\) for every \(g\in G_n\), in particular for \(g\in R\).  The model
built from this \(C\) in the proof of Theorem~\ref{thm:encoding}
satisfies the added constraints.
\end{proof}

For the lift formulas no symmetry breaking is used.

%======================================================================
\section{Certification}
\label{sec:certification}

Every computer-assisted claim below is a claim that an explicit formula
is unsatisfiable.  We produced each formula with the generator scripts in
\path{scripts/} of the accompanying material, recorded its SHA-256 digest,
solved it with CaDiCaL~2.1.2 \cite{BiereEtAl24CaDiCaL} writing a binary
DRAT proof, and checked the proof with \texttt{drat-trim}
\cite{WetzlerHeuleHunt14} at commit \texttt{2e3b2dc0}.  The checker
reported \texttt{s VERIFIED} in every case.  A DRAT proof lists clauses
that are each derivable from the formula and the earlier clauses by a
syntactic rule (reverse unit propagation or resolution asymmetric
tautology); ending in the empty clause, it shows that the formula is
unsatisfiable.  The checker, not the solver, is therefore the component
whose correctness is relied on.

\paragraph{What is trusted.}
The mathematical reductions (Theorem~\ref{thm:encoding},
Propositions~\ref{prop:maximum-encoding}--\ref{prop:lex-sound}) are proved
above.  Their implementation in the generator scripts, the DRAT checker, and
the correspondence between the recorded digests and the checked formulas
are trusted.  As an independent test of the generators, the formulas
without the double-coverage constraint are satisfiable in every case we
tried, and the decoded models were confirmed ample, and in the maximum
case maximum, by a separate direct test.  The encodings also reproduce the
known positive instances: \(L_{B,291}\) for the contraction of Hall's
class at its first coordinate is satisfiable, and its decoded model passes
the direct test of Proposition~\ref{prop:291}.

\paragraph{Archived material.}
All formulas can be regenerated from the scripts, and their digests are
recorded.  For \(F_{n,d}\) with \(n\le6\) and for \(M_{6,4}\), \(M_{7,4}\),
\(M_{7,5}\), \(M_{8,5}\), \(M_{8,6}\), \(F_{7,6}\) and \(F_{7,7}\), the
compressed proofs are archived.  For the remaining cases, namely
\(M_{8,3}\), \(M_{8,4}\), \(F_{7,3}\), \(F_{7,4}\), \(F_{7,5}\) (proofs of
\(0.49\) to \(3.9\)\,GB) and the lift formulas (proofs of about
\(80\)\,MB each), we record the digests of the formula and of the checked
proof and the checker's verdict and statistics.  A separate replay of all
the \(L_{B,290}\) formulas, with retained proofs, is recorded with the
companion work \cite{BousquetDampReader}.  Rechecking them means regenerating the
formula, comparing its digest, running a solver and checking the new proof;
a new proof need not coincide with the original byte for byte.
Table~\ref{tab:certificates} lists the certificates.

%======================================================================
\section{Maximum classes on at most eight coordinates}
\label{sec:maximum}

\begin{proof}[Proof of Theorem~\ref{thm:A}]
Let \(C\) be a maximum class of VC dimension \(d\) on \(n\le8\)
coordinates; all coordinates are active when \(d\ge1\), and \(d=0\) means
a single vertex, which is a corner.  If \(n\le7\), Theorem~\ref{thm:B}
applies.  Let \(n=8\).  By Lemma~\ref{lem:small-vc}, a cornerless
maximum class would have \(3\le d\le6\).  For each of these four values of
\(d\), the formula \(M_{8,d}\) with lex-leader constraints is unsatisfiable
by the certificates in Table~\ref{tab:certificates}, so by
Propositions~\ref{prop:maximum-encoding} and~\ref{prop:lex-sound} no
cornerless maximum class exists.
\end{proof}

The proof of Theorem~\ref{thm:B} below also covers maximum classes on
seven coordinates.  We have additionally refuted \(M_{6,4}\), \(M_{7,4}\) and
\(M_{7,5}\) directly; these certificates are much smaller than those of
Theorem~\ref{thm:B} and are listed for comparison.

\begin{table}[t]
\centering
\small
\begin{tabular}{@{}lrrrrr@{}}
\toprule
Formula & \(\Phi_d(n)\) & variables & clauses & core lemmas & proof\\
\midrule
\(M_{6,4}\) & 57 & 2\,571 & 5\,299 & 69 & archived\\
\(M_{7,4}\) & 99 & 12\,205 & 27\,252 & 20\,269 & archived\\
\(M_{7,5}\) & 120 & 6\,519 & 13\,232 & 140 & archived\\
\(M_{8,3}\) & 93 & 51\,496 & 134\,558 & 2\,383\,367 & 0.88\,GB, digest only\\
\(M_{8,4}\) & 163 & 49\,858 & 118\,878 & 3\,010\,081 & 1.16\,GB, digest only\\
\(M_{8,5}\) & 219 & 34\,752 & 75\,548 & 158\,465 & archived\\
\(M_{8,6}\) & 247 & 16\,386 & 32\,958 & 693 & archived\\
\midrule
\(F_{n,d}\), \(3\le d\le n\le6\) & & \(\le15\,846\) & \(\le39\,835\) & \(\le121\,524\) & archived\\
\(F_{7,3}\) & & 58\,531 & \(\le148\,463\) & 1\,299\,213 & 0.49\,GB, digest only\\
\(F_{7,4}\) & & 58\,531 & \(\le148\,463\) & 7\,972\,306 & 3.90\,GB, digest only\\
\(F_{7,5}\) & & 58\,531 & \(\le148\,463\) & 4\,607\,445 & 2.05\,GB, digest only\\
\(F_{7,6}\) & & 58\,531 & \(\le148\,463\) & 72\,680 & archived\\
\(F_{7,7}\) & & 58\,531 & \(\le148\,463\) & 1 & archived\\
\midrule
\(L_{B,290}\), \(B\) a Hall contraction & & \(\le592\,845\) & \(\le2\,174\,513\) & 78 cases & digest only\\
\(L_{C_H/e,294}\), \(e=8,10\) & & 592\,845 & 2\,174\,513 & \(\le62\,863\) & digest only\\
\bottomrule
\end{tabular}
\caption{Certificates.  All formulas carry lex-leader constraints except the
lift formulas.  ``Core lemmas'' is the number of proof lemmas that
\texttt{drat-trim} retains in its backward check.  Digests, checker logs and
run times are in \texttt{evidence/}.}
\label{tab:certificates}
\end{table}

%======================================================================
\section{Ample classes on at most seven coordinates}
\label{sec:seven}

\begin{proof}[Proof of Theorem~\ref{thm:B}]
Suppose that \(C\) is a nonempty cornerless ample class on at most seven
coordinates.  Discarding inactive coordinates does not change the cube
structure, so we may assume that \(C\subseteq Q_{[n]}\) with \(n\le7\) and
every coordinate active.  Let \(d\) be its VC dimension.  By
Lemma~\ref{lem:small-vc}, \(d\ge3\), so \(3\le d\le n\) and \(n\ge3\).  By
Theorem~\ref{thm:encoding} and Proposition~\ref{prop:lex-sound}, the
formula \(F_{n,d}\) with lex-leader constraints is satisfiable.  But for
each of the \(15\) pairs \(3\le d\le n\le7\) it is unsatisfiable by the
certificates in Table~\ref{tab:certificates}.  This contradiction proves the
theorem.
\end{proof}

\begin{corollary}
\label{cor:seven-peelable}
Every nonempty ample class on at most seven coordinates has a corner
peeling.  Every nonempty cornerless ample class on eight active coordinates
has no peripheral coordinate and is not the product of two classes on
nonempty coordinate sets.
\end{corollary}

\begin{proof}
Deleting a corner from an ample class leaves an ample class
(Theorem~\ref{thm:corner-deletion}) on the same coordinates, which has a
corner again by Theorem~\ref{thm:B} as long as it is nonempty.  For the second
statement, a peripheral coordinate would give a nonempty cornerless ample
class on seven coordinates by Lemma~\ref{lem:peripheral}, and a product
decomposition would give one on at most seven coordinates by
Lemma~\ref{lem:products}; both contradict Theorem~\ref{thm:B}.
\end{proof}

\begin{remark}
\label{rem:cost}
The cost is concentrated in the middle dimensions.  For \(n=7\), the
cases \(d=6,7\) take seconds, whereas \(d=3,4,5\) take between eight
minutes and two hours of solver time and produce proofs of \(0.5\) to
\(3.9\)\,GB.  The extreme dimensions are easy for structural reasons: when
\(d=n\) the complex contains \(E\) and the unique facet cube is the whole
cube, whose vertices are covered once, which the solver sees after
propagation (the backward core has a single lemma).
\end{remark}

%======================================================================
\section{Lifts of Hall contractions}
\label{sec:hall}

\subsection{Hall's class and its contractions}

We use the \(299\) concepts of \(C_H\subseteq Q_{12}\) exactly as listed in
the ancillary file of \cite{ChalopiChepoiMoran22}, with coordinates numbered
\(0,\dots,11\) in the order of that file.  Write \(C_H/e\) for the
contraction at coordinate \(e\), and \(C_H/\{e,f\}\) for the contraction at
two coordinates.

\begin{proposition}
\label{prop:hall-structure}
\begin{enumerate}
\item \(C_H\) is a maximum class of VC dimension \(3\) with
      \(\Phi_3(12)=299\) concepts and no corner.
\item Each single contraction \(C_H/e\) is ample, has order \(232\), uses
      all eleven coordinates, and has exactly two corners.
\item For \(i=0,\dots,5\), the contractions \(C_H/2i\) and \(C_H/(2i+1)\)
      coincide as subsets of \(Q_{11}\), when the remaining coordinates are
      renumbered in increasing order.  The six classes \(C_H/2i\) are
      pairwise distinct.
\item Each double contraction \(C_H/\{e,f\}\) is ample, has order \(176\),
      uses all ten coordinates, and has two, three or four corners.  It has
      exactly two corners if and only if \(\{e,f\}\) is one of the pairs
      \(\{0,1\},\{2,3\},\dots,\{10,11\}\) of item~(3).
\item The only cube automorphism \(g\in G_{12}\) with \(g(C_H)=C_H\) is
      the identity.
\end{enumerate}
\end{proposition}

\begin{proof}
All items are finite checks performed by
\path{scripts/verify_ample_corners_data.py}, which uses no SAT solver.  It
tests ampleness by the Sandwich equality \(|\Sh|=|\SSh|=|X|\) of
Theorem~\ref{thm:sandwich}, and counts corners in two ways: by
Lemma~\ref{lem:facet-cubes}(3), and directly from the definition, by listing
the maximal cubes at each vertex.  The two counts agree in every case.

For (5), write a cube automorphism as \(g(x)=\pi(x)\oplus t\) with \(\pi\) a
coordinate permutation.  Since \(0\in C_H\), the translation \(t=g(0)\)
lies in \(C_H\).  For each of the \(299\) candidates \(t\), the script
assigns \(\pi(0),\pi(1),\dots\) in turn, keeping only partial assignments
that preserve the numbers \(|\{x\in C:x_e=x_f=1\}|\) for all assigned
\(e,f\) (with \(C=C_H\) on one side and \(C_H\oplus t\) on the other).  Each
complete assignment is then tested directly.  Only the identity survives.
\end{proof}

Item~(3) says that the twelve coordinates of \(C_H\) come in six
\emph{twin} pairs: if \(\tau_i\) exchanges \(2i\) and \(2i+1\), then \(C_H\) and
\(\tau_i(C_H)\) have the same contraction at \(2i\), although
\(\tau_i(C_H)\ne C_H\) by item~(5).  Item~(4) detects the same pairs through
corner counts.  Item~(5) means that searches restricted to classes invariant
under a nontrivial group of cube symmetries cannot find Hall's class.

\subsection{Optimal lifts}

\begin{proposition}
\label{prop:291}
For \(e\in\{0,2,4,6\}\) there is a cornerless ample class
\(Y_e\subseteq Q_{12}\) of order \(291\), and for \(e\in\{8,10\}\) one of order
\(295\), such that all coordinates of \(Y_e\) are active and contracting its
last coordinate gives \(C_H/e\).  Each \(Y_e\) has VC dimension \(3\).
Consequently \(\cmin=\mamp\le291\).
\end{proposition}

\begin{proof}
The classes are listed in \path{evidence/hall_lifts/}.  They were obtained
by decoding satisfying assignments of \(L_{C_H/e,\,291}\) and
\(L_{C_H/e,\,295}\), but their properties are checked directly by
\path{scripts/verify_ample_corners_data.py}: \(|\Sh(Y_e)|=|\SSh(Y_e)|=|Y_e|\),
no vertex lies in a unique maximal cube, all coordinates are active, and
contracting the last coordinate gives \(C_H/e\).  Of the \(232\) vertices of
\(C_H/e\), exactly \(|Y_e|-232\) lift to both layers.  The last statement
follows from Proposition~\ref{prop:cmin-equals-mamp}.
\end{proof}

The companion work constructs a class of order \(291\) over \(C_H/0\) by
hand, starting from an ample overlap of order \(59\)
\cite{BousquetDampReader}.

\begin{proof}[Proof of Theorem~\ref{thm:C}]
Let \(B\) be \(C_H/e\) or \(C_H/\{e,f\}\), on the coordinate set \(E\), and let
\(Y\) be a cornerless ample lift of \(B\) with \(|Y|\le N\).  Every
coordinate of \(E\) is active in \(Y\), because \(\pi_zY=B\) and all
coordinates of \(B\) are active by Proposition~\ref{prop:hall-structure}.
The coordinate \(z\) is active as well.  Otherwise \(Y=B\times\{i\}\) is a
copy of \(B\) in one layer, which has the same cubes and corners as \(B\);
but \(B\) has a corner by Proposition~\ref{prop:hall-structure}.  Hence
Proposition~\ref{prop:lift-encoding} shows that \(L_{B,N}\) is satisfiable.

Take \(N=290\).  For each of the \(12\) single and \(66\) double contractions,
\(L_{B,290}\) is unsatisfiable by the certificates listed in
\path{evidence/hall_lifts/certificates.txt}.  This proves \(|Y|\ge291\).

For single contractions, the upper bounds are
Proposition~\ref{prop:291} for even \(e\), and the same classes serve for
odd \(e\), because \(C_H/(2i+1)=C_H/2i\) by
Proposition~\ref{prop:hall-structure}(3).  For \(e\le7\) this matches the
lower bound \(291\).  For \(e\in\{8,10\}\), the formula \(L_{C_H/e,\,294}\) is
unsatisfiable by the certificates in \path{evidence/hall_lifts/order294/};
by the argument above, no cornerless ample lift of \(C_H/e\) has order at
most \(294\).  The same holds for \(e\in\{9,11\}\), whose base is the
same class.  With Proposition~\ref{prop:291}, the optimum for \(e\ge8\) is
\(295\).
\end{proof}

\begin{remark}
\label{rem:hall-doubles}
For double contractions we only have the lower bound \(291\).  The class
\(C_H/e\) is itself a lift of \(C_H/\{e,f\}\), but it has two corners, so it
gives no upper bound.  The dichotomy between the twin pairs
\(\{0,1\},\dots,\{6,7\}\) and \(\{8,9\},\{10,11\}\) in Theorem~\ref{thm:C} is
not explained by any symmetry of \(C_H\), which has none.
\end{remark}

%======================================================================
\section{Contract-and-relift descent}
\label{sec:descent}

Theorem~\ref{thm:C} describes one step of a procedure that can be
iterated.  Let \(Y\) be a cornerless ample class.  Contract it at a coordinate
\(c\), obtaining an ample class \(B=\pi_cY\) (which, in every case below, has
corners), and search for a cornerless ample lift of \(B\) of least order with
the formula \(L_{B,N}\) of Section~\ref{sec:lift-formula}.  The lift need not
contain \(Y\) or be isomorphic to it.  Starting from Hall's class, the first
step gives the classes \(Y_e\) of Proposition~\ref{prop:291}.  We iterate
from these.

\begin{proof}[Proof of Theorem~\ref{thm:D}]
The class is listed in \path{evidence/descent/record.json}.  The script
\path{scripts/verify_ample_corners_data.py} checks directly, without a SAT
solver, that it has \(267\) concepts on twelve active coordinates, that
\(|\Sh|=|\SSh|=267\), so that it is ample by Theorem~\ref{thm:sandwich}, that
its VC dimension is \(3\), and that no vertex lies in a unique maximal cube
(by both corner tests).  The equality \(\cmin=\mamp\) is
Proposition~\ref{prop:cmin-equals-mamp}, and \(\cmin\ge64\) is proved in
\cite{BousquetDampReader}.
\end{proof}

The record was found as follows; only the final class matters for the
theorem, but the lineage shows how the method moves away from Hall's
class.  Each arrow is ``contract at the indicated coordinate, then take the
decoded lift of least order found''.
\[
 C_H\ (299)\ \xrightarrow{\,0\,}\ Y_0\ (291)\ \xrightarrow{\,0\,}\ 283\
 \xrightarrow{\,0\,}\ 275\ \xrightarrow{\,0\,}\ 267 .
\]
In each step the solver was run with bound one below the current record
and, after each satisfying assignment, with bound one below the order found.
The script \path{scripts/descent.py} automates this and verifies every
decoded class directly before storing it; all \(23\) classes it found are in
\path{evidence/descent/state.json}, with their parents, and are rechecked,
including the contraction relation to their parents, by
\path{scripts/verify_ample_corners_data.py}.

\subsection{The shattered sets of the record}

The class of Theorem~\ref{thm:D} has a transparent shattering complex.
Label its coordinates \(0,\dots,11\) as in
\path{evidence/descent/record.json}, and put
\[
 \gamma(i)=\lfloor i/4\rfloor\in\{0,1,2\},\qquad
 \lambda(i)=\lfloor (i\bmod 4)/2\rfloor\in\{0,1\}.
\]
Thus the coordinates fall into three \emph{groups}
\(\{0,1,2,3\},\{4,5,6,7\},\{8,9,10,11\}\), and each group into two
\emph{pairs} \(\{4g,4g+1\}\) (label \(0\)) and \(\{4g+2,4g+3\}\) (label \(1\)).
Call a triple \(\{i,j,k\}\) an \emph{even transversal} if it meets every
group once and \(\lambda(i)+\lambda(j)+\lambda(k)\) is even.  There are
\(4\cdot2^3=32\) even transversals.

\begin{proposition}
\label{prop:record-complex}
The class \(Y_{267}\) of Theorem~\ref{thm:D} shatters exactly the empty set,
the \(12\) singletons, all \(66\) pairs, and the \(188\) triples that are not
even transversals.  Consequently every contraction \(\pi_eY_{267}\) has
\(208\) concepts, the facet cubes of \(Y_{267}\) are its \(188\) triple
cubes, and its only cube automorphism is the identity.
\end{proposition}

\begin{proof}
The shattered sets, the orders of the contractions and the automorphism
group are computed directly by \path{scripts/verify_ample_corners_data.py}
and \path{explore_ample_267_structure.py}.  The count
\(1+12+66+188=267\) agrees with \(|Y_{267}|\) by the Sandwich equality.
Since \(Y_{267}\) is ample, \(|\pi_eY_{267}|=|Y_{267}|-|Y_{267}^{\{e\}}|\) and
\(|Y_{267}^{\{e\}}|\) is the number of shattered sets containing \(e\), namely
\(1+11+47=59\) for every \(e\): each coordinate lies in \(8\) even
transversals.  Every pair lies in at least \(10-2=8\) shattered triples, so no
pair is a facet, and the facets are the \(188\) shattered triples
(Lemma~\ref{lem:facet-cubes}).
\end{proof}

In the complex of Proposition~\ref{prop:record-complex}, every pair of
coordinates lies in either \(0\) or exactly \(2\) even transversals, and the
pairs lying in none form three disjoint copies of \(K_4\), one on each group.
All \(16\) classes of order \(267\) found by the descent, which fall into
\(13\) isomorphism types, have a complex isomorphic to this one, although
their facet-cube placements differ and none has a nontrivial cube
automorphism.  The same twin pairs appear in Hall's class
(Proposition~\ref{prop:hall-structure}): the descent leading to \(Y_{267}\)
contracts coordinate \(0\) at every step and so only rotates Hall's
coordinate labels.  Nevertheless \(Y_{267}\) is far from Hall's class.
Among the copies of \(C_H\) obtained by a cyclic rotation of the coordinates
followed by a translation, the best one shares about \(100\) of the \(267\)
vertices and \(36\) of the \(188\) facet cubes.

\begin{remark}[Local optimality, not certified]
\label{rem:descent-local}
Starting the descent from each of the other optimal lifts \(Y_2,Y_4,Y_6\)
(order \(291\)) and \(Y_8,Y_{10}\) (order \(295\)) also ends at exactly
\(267\).  Several further searches for a cornerless ample class of order at
most \(266\) returned no solution: one-coordinate lifts of all single and
double contractions of the order-\(267\) classes, including all \(556\)
double contractions, whose lifts have eleven coordinates; neighbourhood
searches that keep the facet cubes of an order-\(267\) class on all triples
avoiding a chosen set of free coordinates and allow at most \(187\)
triples; and prescribed complexes with the same parity structure on nine to
eleven coordinates.  In the neighbourhood searches, the same formulas with
\(188\) triples allowed are satisfiable within seconds, and their decoded
models are cornerless ample classes of order \(267\) that differ from the
source class, so the free part is genuinely free.  These runs did not
record proofs, and we make no claim based on them.  They indicate that
\(267\) is a strong local minimum of the moves we have tried.
\end{remark}

%======================================================================
\subsection{Certified shrinking edits of the record}
\label{sec:record-shrinking-edits}

We can exclude two small neighborhoods of the particular labelled class
\(X=Y_{267}\) by complete enumeration, including full peeling certificates
for every ample output.  The intermediate deletion need not be ample.

\begin{proposition}[Computational exclusion of small shrinking edits]
\label{prop:record-small-shrinking-edits}
Let \(R\subseteq X\), let \(A\subseteq Q_{12}\setminus X\), and suppose
\(|R|=r\), \(|A|=r-1\), where \(r\in\{2,3\}\).
If \(Y=(X\setminus R)\cup A\) is ample, then \(Y\in\Damp\).
There are exactly \(24\) such labelled ample outputs for \(r=2\), and
exactly \(1296\) for \(r=3\).
\end{proposition}

\begin{proof}[Enumeration and certificate verification]
Every output has \(266\) vertices, so ampleness requires
\(|\Sh(Y)|=266\).  By monotonicity of shattering, at least one of the
\(267\) supports in \(\Sh(X)\) must cease to be shattered by
\(X\setminus R\).  Such a support \(S\) is lost precisely when \(R\)
contains a whole nonempty trace fiber
\[
 F_{S,t}=\{x\in X:x|_S=t\},\qquad t\in\{0,1\}^S.
\]
The verifier enumerates these fibers for every \(S\in\Sh(X)\).
There are no singleton fibers.  There are \(24\) distinct two-element
fibers and \(6\) distinct three-element fibers.  Consequently only
\(24\) deletion pairs and \(6354\) deletion triples need be considered;
all other choices retain at least \(267\) shattered supports.

For a remaining deletion set put \(B=X\setminus R\), and, for every
coordinate support \(S\), let
\(M_B(S)=\{0,1\}^S\setminus\{x|_S:x\in B\}\).
The exact completion test is
\[
 S\in\Sh(B\cup A)
 \quad\Longleftrightarrow\quad
 M_B(S)\subseteq\{a|_S:a\in A\}.
\]
In particular, supports missing more than \(|A|\) traces can be ignored.
Every inserted word must individually create at most
\(266-|\Sh(B)|\) supports.  This necessary filter leaves \(8652\)
pairs of inserted words in total in the three-deletion case; the joint
completion test then gives exactly the \(1296\) ample outputs stated.
For two deletions, each of the \(24\) pairs admits exactly one insertion.

The archives \path{evidence/order267_two_for_one/} and
\path{evidence/order267_three_for_two/} contain all resulting classes and
complete \(266\)-vertex corner orders.  Their corresponding
\path{scripts/verify_order267_*.py} verifiers check coverage, recompute
shattering, and check that the full cube on the current incident
directions of each deleted vertex is present.  The three-deletion verifier
recomputes every candidate's projections directly from its membership
bitset, independently of the exploration code's trace-update rule.
Every order ends with the empty class, proving the assertion.
\end{proof}

This excludes the stated edits of one seed.  It neither excludes larger
edits nor bounds the size of an arbitrary cornerless ample class.

\begin{proposition}[Corner-deleted contractions have no cornerless ample lift]
\label{prop:record-corner-deleted-relifts}
Let $X=Y_{267}$, let $e$ be any of its twelve coordinates, and let $c$
be any corner of $\pi_eX$. Put $B=(\pi_eX)\setminus\{c\}$.
There is no cornerless ample class $Y\subseteq B\times\{0,1\}$ whose
contraction of the last coordinate is exactly $B$.
There is no restriction on the order or VC dimension of $Y$.
\end{proposition}

\begin{proof}[Computer-assisted proof]
After suppressing coordinate $e$, encode the remaining eleven coordinates
by their binary integer values, starting at coordinate zero.
Each contraction has $208$ vertices and exactly two corners:
\[
\begin{array}{c|c}
e&\text{corners of }\pi_eX\\\hline
0,1&127,217\\
2,3&1287,1860\\
4,5&283,366\\
6,7&1264,1536\\
8&236,1736\\
9&492,1992\\
10,11&1330,1571
\end{array}
\]
These lists are checked directly from the incident cubes. Deleting a
corner leaves an ample $207$-vertex base. Within each paired row of the
table the bases are identical. The two bases for $e=9$ are translates
of the corresponding bases for $e=8$, by binary exclusive-or with $256$.
Thus twelve explicit base families cover all twenty-four choices $(e,c)$.

For each representative, use a membership variable $p_v$ for every
$v\in E=B\times\{0,1\}$. Apply the exact ampleness and noncorner
conditions from the proof of
Theorem~\ref{thm:record267-ample-inclusion-minimality}, with potential
family $E$. Add precisely the contraction clauses
\[
 p_{(x,0)}\lor p_{(x,1)}\qquad(x\in B).
\]
Membership outside $E$ is forbidden by the choice of variables. These
clauses therefore impose $\pi_zY=B$, rather than merely containment.
The other clauses impose ampleness and absence of corners exactly.
No cardinality, prescribed-shadow, or VC-dimension clauses are added.
Consequently a model would be exactly a class forbidden by the statement.
All twelve formulas have DRAT refutations checked independently by
\texttt{drat-trim}.

The formulas, compressed proofs, hashes, coverage translations and checker
logs are in \path{evidence/corner_deleted_relifts_unbounded/}.
The generator is \path{scripts/explore_corner_deleted_relifts.py};
\path{scripts/verify_corner_deleted_relifts.py} reconstructs every formula,
checks all twenty-four input identifications, and replays the proofs.
It also exhausts the contraction module on three-coordinate families and
checks that the unchanged contraction of $X$ admits the known lift $X$.
\end{proof}

This excludes changing the contraction by one corner deletion and then
choosing an entirely new lift, even one using words absent from $X$.
It does not exclude lifts of other bases, further deletions from these
bases, or cornered nonpeelable lifts. In particular, it does not improve
the unrestricted interval $64\le c_{\min}\le267$.

The exclusion cannot be explained by incompatible choices at the base
corners alone. For each of the twelve representative bases, the archive
\path{evidence/relift_corner_localization/} gives consistent layer
memberships on the union of the incident corner cubes. On each such cube
the induced lift is ample and neither lift of the base corner is a
corner. These are local assignments, not full ample lifts.
More strongly, for $B=(\pi_0Y_{267})\setminus\{127\}$,
\path{evidence/relift_ampleness_rank_isometric/report.json} records a
$272$-vertex isometric cornerless lift with contraction $B$ such that every shattered support of
size at most three is strongly shattered. Direct counting gives $273$
shattered supports and exactly three that are not strongly shattered;
all three have size four. Thus even imposing isometry, cornerlessness
globally, and all ampleness equalities through rank three does not suffice
for this base. Each of the $142$ possible single-word insertions preserving
the contraction creates a new shattered support, so none repairs ampleness.
These nonample diagnostics do not change the minimum-order bounds.

\begin{proposition}[All ample subfamilies of the $272$-vertex diagnostic peel]
\label{prop:native272-cycle-fiber}
Let $T\subseteq Q_{12}$ be the $272$-vertex isometric family just
described, with coordinates numbered from zero, and put
\[
 U=\{2150,2166,2174,2278,2286\},\qquad
 F=\{x\in T:x_8=0,\ x_{11}=x_3=x_4=x_7=1\}.
\]
Then $|F|=17$, and the following hold.
\begin{enumerate}
\item Every ample $Y\subseteq T$ satisfies
      $U\nsubseteq Y$ or $Y\cap F=\varnothing$.
\item Every ample subfamily of $T$ belongs to $\Damp$.
\end{enumerate}
Thus no deletion of an arbitrary batch of words from this fixed family
produces an ample forbidden pc-minor, of any order.
\end{proposition}
\begin{proof}[Restriction argument and computer-assisted exclusion]
Restrict $T$ to $x_8=0,x_{11}=1$ and project onto coordinates $(3,4,7)$.
The resulting three-coordinate family is $Q_3\setminus\{100,011\}$,
where the omitted words are antipodal; equivalently, its integer encoding
is $\{0,2,3,4,5,7\}$, with the first coordinate the least significant bit.
It is a six-cycle, with six vertices and seven shattered supports, so it
is not ample. Five trace fibers are precisely the singletons indexed by
$U$; the sixth is $F$. These image and fiber identities are checked
directly from the recorded words. If $Y$ retained all five singleton
fibers and any word of $F$, its same restriction and projection would
still be this six-cycle. Closure of ample classes under restriction and
projection proves the first assertion.

First exclude cornerless ample subfamilies avoiding $F$. Encode arbitrary nonempty subfamilies of $T$
using the exact ampleness and noncorner conditions of
Theorem~\ref{thm:record267-ample-inclusion-minimality}. Add the just-proved
condition that at least one of the six trace fibers is empty, and force
every membership variable in $F$ to be false. The last condition already
satisfies the additional trace condition, so the query covers every
nonempty ample cornerless subfamily of $T\setminus F$. There is no order
bound or restriction on which other vertices may be removed. The formula
has $18166$ variables and $155690$ clauses and has an independently
replayed DRAT refutation. Hence every nonempty ample subfamily of
$T\setminus F$ has a corner. Corner deletion preserves ampleness and
containment in $T\setminus F$, so induction gives a complete peeling.

This first formula, compressed proof, and replay log are retained in
\path{evidence/native272_subfamilies_unbounded_cycle_cut_trace152/}.
The verifier \path{scripts/verify_native272_cycle_cut.py} checks the
six-cycle image and all fibers, reconstructs the formula, and replays
the refutation.

It remains to consider a nonempty cornerless ample $Y\subseteq T$
meeting $F$. List $U$ in the displayed increasing order as
$u_0,\ldots,u_4$. The first assertion gives a least index $j$ with
$u_j\notin Y$. There are exactly five disjoint cases, given by
\[
 u_0,\ldots,u_{j-1}\in Y,\qquad u_j\notin Y,\qquad
 Y\cap F\ne\varnothing \qquad(0\le j\le4).
\]
For each case add these membership conditions to the same exact ample
cornerless-subfamily encoding, without a cardinality or VC bound.
Each formula has $18160$ variables; its clause count is $155646+j$.
All five have DRAT refutations, first checked on the search worker and
then independently replayed locally after reconstructing the formulas
byte for byte. The formulas, compressed proofs, hashes, and logs are in
\path{evidence/native272_remote_split/}; the reconstruction and replay
are in \path{scripts/verify_native272_remote_split.py}.
This excludes every nonempty cornerless ample subfamily of $T$.
Any nonempty ample $Y\subseteq T$ therefore has a corner; its deletion
leaves another ample subfamily of $T$. Induction on $|Y|$ gives a full
corner peeling and proves the second assertion. The conclusion concerns
this fixed envelope, not arbitrary classes of order at most $272$.
\end{proof}

\begin{proposition}[Relative minimality of the $264$-word rooted input]
\label{prop:root264-relative-ample-inclusion}
Let $A=(Y_{267}\setminus\{236,492,1004,748\})\cup\{220\}$ be the
recorded $264$-word rooted input, with marked root $r=750$.
Every proper ample $B\subsetneq A$ peels relatively to $B\cap\{r\}$.
If $r\notin B$, this means an ordinary peeling to the empty family.
In fact every maximal corner-deletion sequence avoiding $r$ reaches
that target. This assertion concerns the specified root, not arbitrary
rerootings of proper subfamilies.
\end{proposition}
\begin{proof}[Computer-assisted proof]
Use the membership encoding of
Theorem~\ref{thm:record267-ample-inclusion-minimality}, with potential
family $A$. Omit the noncorner condition at $r$, leave its membership
variable free, require at least one vertex outside $\{r\}$, and require
a proper subfamily. Thus a model is exactly a proper ample $B\subset A$
with $B\setminus\{r\}\ne\varnothing$ and no corner outside $\{r\}$.
The resulting formula has $17008$ variables and $145304$ clauses.
It has an independently checked DRAT refutation. Therefore each proper
ample subfamily not already contained in $\{r\}$ has a deletable
corner outside $\{r\}$. Deleting such a corner preserves ampleness,
proper containment in $A$, and the intersection with $\{r\}$.
Induction proves the assertion, regardless of deletion choices.

The script \path{scripts/explore_root264_inclusion.py} constructs the
formula. The proof, hashes and checker log are in
\path{evidence/root264_ample_inclusion/}. The checker
\path{scripts/verify_root264_inclusion.py} rebuilds that formula,
verifies the input's ordinary corner order and unique corner, tests the
protected-root encoding in $768$ small membership cases, and optionally
replays the compressed proof with \texttt{--replay-proof}.
The unchanged ordinary encoding is separately rebuilt against its
previously checked $267$ certificate.
\end{proof}

The stronger endpoint statement remains open: does every proper ample
$B\subsetneq A$ admit a corner order ending at each prescribed $b\in B$?
By Proposition~\ref{prop:root264-relative-ample-inclusion}, every such $B$
is ordinarily positive. Consequently the stronger statement is equivalent
to the absence of a proper ample subfamily with at least two vertices and
exactly one corner. In one direction that unique corner cannot be retained
until the end. In the other, a failed maximal peeling avoiding a prescribed
vertex leaves a nontrivial ample remainder with no corner outside that
vertex. The remainder is proper in $A$ and ordinarily positive, so its
prescribed vertex is its unique corner. This equivalence reduces the
endpoint question to a single exact membership search with a free choice
of root; it does not settle that search. The formulations and their current
unresolved status are retained in
\path{evidence/root264_free_root/} and
\path{evidence/root264_free_root_exact/}.

The vertex-gluing criterion in the companion Damp manuscript now shows
that the $527$-word double has every proper ample subfamily in $\Damp$.
Joining the two roots by a new edge instead gives a $528$-word cornerless
class with the same all-subfamily property. It is not a forbidden pc-minor,
since the bridge contracts to the negative $527$ class. The uniform
criterion and full argument appear there as
\texttt{prop:damp-inclusion-minimal-cut-vertex} and
\texttt{cor:damp-inclusion-minimal-527-528}. This distinguishes minimality
under ample inclusion from the contraction minimality required in the
forbidden-minor problem; it gives no improvement to the global size bounds.

\begin{remark}[Batch deletion in two-row presentations]
\label{rem:record-two-row-fiber-deletions}
The labelled $Y_{267}$ has six unordered coordinate pairs for which,
after a sign choice, its four rows have the form
$(B_0,D_0,D_1,B_1)$ with $B_0\cup B_1=D_0\cap D_1=:C$.
One pair has row sizes $13,136,72,46$ and $|C|=48$.
For each presentation and either large row, delete a nonempty fiber
$\{x\in D_i:x|_S=t\}$ disjoint from $C$, leaving all other rows fixed.
All $7148$ distinct resulting families are nonample; the largest deletion
has $56$ words. This extends the preceding tests to a particular family
of larger batches without insertions, not to arbitrary larger edits.

The independent support/trace enumeration in the companion script
\path{../dismantlable-ample/research/verify_two_row.py} agrees with the
explorer's iterated bitset-intersection enumeration. For every output it
records a shattered support with no witnessing cube, certifying failure
of ampleness. A separate shadow-cardinality computation checks every
output, including those with corners. Certificates are in
\path{../dismantlable-ample/research/two_row_verification.json}.
No change to the $64$--$267$ interval follows.
\end{remark}

\begin{proposition}[Rigidity with a fixed two-row contraction]
\label{prop:record-fixed-contraction-rigidity}
Fix any one of the six two-row presentations in
Remark~\ref{rem:record-two-row-fiber-deletions}, and retain its large
rows $D_0,D_1$. Put $C=D_0\cap D_1$. For arbitrary subsets
$B_0,B_1\subseteq C$ with $B_0\cup B_1=C$, define
\[
 V=(B_0\times\{00\})\cup(D_0\times\{01\})
   \cup(D_1\times\{10\})\cup(B_1\times\{11\}).
\]
If $V$ is ample and cornerless, then $(B_0,B_1)$ is either the original
pair of diagonal rows or its exchange. Thus there are exactly two
labelled assignments for each fixed pair $D_0,D_1$, both of order $267$.
No bound on $|B_0\cap B_1|$ is assumed.
\end{proposition}
\begin{proof}[Exact encoding and certificate verification]
Each fixed pair satisfies $|D_0|+|D_1|=208$ and $|C|=48$.
The family $K=(D_1\times\{0\})\cup(D_0\times\{1\})$ is ample.
By the two-row ampleness identity in the companion Damp manuscript,
$V$ is ample exactly when $H=(B_0,B_1)$ is ample. We describe an exact
finite encoding rather than testing a neighbourhood of the original rows.

For each $x\in C$ introduce membership variables $a_{x,0},a_{x,1}$,
with clause $a_{x,0}\vee a_{x,1}$. These allow all three possible
memberships of each word. Every support of $H$ avoiding its new coordinate
is in $\Sh(C)$, independently of the assignment. For every $S\in\Sh(C)$,
require that some $S$-cube contained in $C$ lie wholly in one of the two
sections. This is exactly strong shattering of $S$ in $H$.
The support $S\cup\{z\}$ is shattered precisely when both sections
contain every $S$-trace. Whenever this occurs, require that some
$S$-cube of $C$ lie wholly in $B_0\cap B_1$. This is exactly strong
shattering of $S\cup\{z\}$. Supports outside $\Sh(C)$ cannot be
shattered even after adding $z$. These requirements therefore express
$\Sh(H)=\operatorname{sSh}(H)$ exactly.

For completeness, cornerlessness is encoded directly on the fixed
potential family
\[
 E=(C\times\{00,11\})\cup(D_0\times\{01\})
   \cup(D_1\times\{10\}).
\]
Write $p_v$ for the membership variable of a diagonal vertex, for the
constant $1$ on the two retained rows, and for $0$ outside $E$.
Let $N_E(v)$ be the set of coordinate directions $i$ with
$v^{\{i\}}\in E$. For every $v\in E$ impose
\[
 p_v\ \Longrightarrow\
 \bigvee_{\substack{T\subseteq N_E(v)\\ |T|\ge2}}
 \left[\left(\bigwedge_{i\in T}p_{v^{\{i\}}}\right)
 \wedge\neg p_{v^T}\right].
\]
The disjunction supplies a missing vertex of the cube generated by
present neighbours. It is equivalent to saying that a present $v$ is
not a corner: a missing cube vertex cannot differ in zero or one
incident direction. Thus no degree surrogate or assumed VC bound enters.
Standard conjunction/disjunction gate clauses convert these exact
conditions to a conjunctive normal form, preserving their models.

For each of the six fixed contractions, complete model enumeration gives
exactly two assignments to the membership variables. Each is checked
directly for ampleness and cornerlessness and identified with the original
rows or their exchange. A clause excluding each of these two assignments
is then added. The resulting formula has an independently checked DRAT
refutation. Since every other assignment satisfies both exclusion clauses,
this proves completeness. In particular the search has no cardinality
cutoff. Separate checked refutations with $|B_0\cap B_1|\le10$ agree
with the resulting lower bound $|V|\ge208+48+11=267$.

The scripts in the companion directory
\path{../dismantlable-ample/research/} are
\path{two_row_fixed_contraction_sat.py}, \path{two_row_fixed_rigidity.py},
and \path{verify_two_row_fixed_contraction.py}. Their
\path{two_row_fixed_contraction/} archive contains the formulas, two
surviving assignments per pair, refutations, checker logs, and hashes.
The replay reconstructs each fixed pair from the original labelled seed,
reconstructs the blocking clauses from the recorded assignments, checks
the resulting formula, and independently replays each refutation.
\end{proof}

This excludes every change of the diagonal rows with those large rows
fixed, not merely the earlier cylinder deletions. It is a finite rigidity
statement for six specified contractions, not a global $267$ lower bound.
Any smaller cornerless construction through these presentations must
change at least one large row. Cornered negative assemblies are not
classified by this test.

\begin{proposition}[A two-row overlap with no negative ample lift]
\label{prop:record-two-row-overlap-envelope}
Use the labelled $Y_{267}$ and the two-row presentation on coordinates
$0,1$, with coordinate $1$ flipped. Put $C=D_0\cap D_1$, deleting
these two coordinates, and $E=C\times Q_1$. Then $|C|=48$, $|E|=96$,
and every ample subfamily of $E$ belongs to $\Damp$.
In particular, no ample forbidden class $H=(B_0,B_1)$ with
$B_0\cup B_1=C$ can serve as the reduced input of a two-row assembly
with this common row intersection. This holds even if the large rows
are changed while their intersection remains this labelled $C$.
\end{proposition}
\begin{proof}[Computer-assisted exclusion and induction]
In integer coordinates, with coordinate zero least significant, the
base is recovered directly from the seed by
\[
 C=\{\lfloor x/4\rfloor:x\in Y_{267},\ x\mathbin{\mathrm{xor}}1\in Y_{267}\}.
\]
Thus it is the contraction in coordinate $1$ of the reduction in
coordinate $0$, with both coordinates removed. The certificate verifies
that $C$ is ample and has $48$ vertices.

Apply the exact all-subfamily encoding used in
Theorem~\ref{thm:record267-ample-inclusion-minimality} to the potential
family $E$, without a properness clause. All $96$ membership variables
are free. Require nonemptiness, ampleness by shattered/strongly-shattered
support equality, and cornerlessness by a missing vertex of the full
incident cube at each retained vertex. No projection, cardinality, VC,
or coordinate-support constraint is imposed. The resulting formula has
$3222$ variables and $21485$ clauses. Its DRAT refutation is checked,
and an independent replay reconstructs the base from the displayed
edge formula and rebuilds the CNF before checking the proof again.
Consequently every nonempty ample subfamily of $E$ has a corner.
Deleting a corner preserves ampleness and containment in $E$; induction
on order proves dismantlability. Any proposed $H$ in the statement is
an ample subfamily of $E$, proving the last assertion.

The complete formula, compressed proof, hashes and two checker logs
are in \path{evidence/two_row_overlap_envelope/}; the construction and
replay scripts are \path{scripts/explore_two_row_overlap_envelope.py}
and \path{scripts/verify_two_row_overlap_envelope.py}.
\end{proof}

Unlike fixed-contraction rigidity, this checks cornered as well as
cornerless reduced inputs, and allows their projection to be a proper
subfamily of $C$. It excludes one specified $96$-vertex envelope,
not all ample classes of order at most $96$. It neither establishes
a global lower bound nor resolves the inverse construction with an
arbitrary forbidden reduced input.

\begin{theorem}[Every proper ample subfamily of the recorded seed peels]
\label{thm:record267-ample-inclusion-minimality}
Let $X=Y_{267}$ be the labelled family in the certificate of order $267$.
Every proper ample subfamily $Y\subsetneq X$ belongs to $\Damp$.
Indeed, any maximal sequence of corner deletions from such a $Y$ ends
at the empty family. In particular $X$ has no nonempty proper ample
cornerless subfamily.
\end{theorem}
\begin{proof}[Computer-assisted exclusion and induction]
We first exclude every nonempty proper ample cornerless $Y\subset X$,
without fixing any row, coordinate support, or deletion radius.
Introduce a membership variable $p_x$ for each $x\in X$, and clauses
\[
 \bigvee_{x\in X}p_x,\qquad \bigvee_{x\in X}\neg p_x,
\]
requiring $Y$ to be nonempty and proper. For every $S\in\Sh(X)$,
let $\mathcal Q_S(X)$ be the family of all axis-parallel $S$-cubes
contained in $X$. Impose
\[
 \left[\bigwedge_{t\in\{0,1\}^S}
       \bigvee_{\substack{x\in X\\x|_S=t}}p_x\right]
 \ \Longrightarrow\
 \left[\bigvee_{Q\in\mathcal Q_S(X)}\ \bigwedge_{x\in Q}p_x\right].
\]
The antecedent means that $Y$ shatters $S$ and the consequent that it
strongly shatters $S$. Every support shattered by $Y$ already belongs
to $\Sh(X)$, so these implications express ampleness exactly. In
particular, losing shattered supports is allowed; none of the original
supports is prescribed to remain.

Use the exact missing-incident-cube conditions from the proof of
Proposition~\ref{prop:record-fixed-contraction-rigidity}, now with the
potential family $E=X$ and every membership variable free. They express
cornerlessness exactly. Gate clauses yield a formula with $17520$
variables and $149656$ clauses. The full seed satisfies the same formula
before the properness clause is added. After that clause is added,
the saved DRAT refutation is independently verified. Hence no nonempty
proper ample cornerless subfamily exists.

Now let $Y\subsetneq X$ be ample. If $Y$ is nonempty, the exclusion
just proved supplies a corner. Removing it preserves ampleness and
leaves a proper subfamily of $X$. Repeat. Cardinality strictly decreases,
and every nonempty remainder still has a corner, so the process ends
at the empty family regardless of the choices. This proves the theorem.

The encoder is \path{scripts/explore_seed267_subfamilies.py}; its
saved formula and compressed refutation are in
\path{evidence/seed267_all_subfamilies/}. The encoding audit
\path{scripts/verify_seed267_subfamily_encoding.py} reconstructs the
saved formula and checks ampleness and cornerlessness separately and
jointly on all $256$ labelled subfamilies of $Q_3$, for $768$ checks.
These small checks supplement the symbolic coverage argument; they do
not replace the full refutation. The independent replay is
\path{scripts/replay_seed267_subfamilies.py}, with formula, proof, seed,
and checker hashes recorded in the evidence directory.
\end{proof}

This rules out all deletion-only improvements of this particular seed,
including smaller cornered obstructions. Any smaller negative ample class
obtained by replacing vertices of this seed must introduce a word outside
it. The statement is additional to pc-minor minimality: it concerns
arbitrary ample subfamilies, rather than coordinate restrictions and
contractions. It neither raises the global lower bound $64$ nor excludes
smaller examples with different vertex sets.

%======================================================================
\subsection{Smaller rooted inputs from the shrinking outputs}
\label{sec:shrinking-rooted-inputs}

The positive outputs of Proposition~\ref{prop:record-small-shrinking-edits}
also supply inputs to the cut-vertex construction of
\cite{BousquetDampReader}.  Here a rooted input is a positive ample family
with a distinguished vertex that cannot be the last vertex of a corner
peeling, although every proper root-preserving pc-minor can peel to the
image of that vertex. Root-preserving minors use contractions and
restrictions containing the root.

\begin{proposition}[A $527$-vertex cut-vertex obstruction]
\label{prop:rooted264-double527}
In the integer encoding of $X=Y_{267}$, put
\[
 A=(X\setminus\{236,492,1004,748\})\cup\{220\},\qquad a=750.
\]
Then $(A,a)$ is a rooted input of order $264$, and every vertex of every
nonempty proper pc-minor of $A$ is attainable as a peeling endpoint.
Gluing two copies at their
roots, on disjoint coordinate blocks, gives a cornerless ample forbidden
pc-minor of order $527$ and isometric dimension $24$.
\end{proposition}
\begin{proof}
The ample $266$-vertex output
$(X\setminus\{236,492\})\cup\{220\}$ has forced first vertices
$1004,748,750$. Deleting its first two corners leaves $A$, which is
ample, peelable, and has unique corner $a$. Thus it cannot peel to $a$.
The archive \path{evidence/shrinking_rooted_inputs/} supplies complete
root-ending orders for all $24$ elementary root-preserving minors of
$A$, checked by \path{scripts/verify_shrinking_rooted_inputs.py}.
Peeling to the root is preserved under subsequent root-preserving minors,
so these orders prove rooted minimality. In fact the stronger certificate
\path{evidence/root264_minor_endpoints/} gives all $5637$ endpoint orders
across all $36$ elementary proper pc-minors of this $A$. Its independent
verifier reconstructs every minor and checks the USO out-map of every order.
For any further minor and desired endpoint, choose a preimage vertex;
rooted minor closure transports the corresponding order. This proves
that no proper pc-minor can supply a rooted input at any new root.

Translate $a$ to zero and put two copies on disjoint coordinate blocks.
Their union is ample: no mixed support is shattered, and the two shadows
share only the empty support, giving $264+264-1=527$ supports and vertices.
Every nonroot vertex remains a noncorner, while the root has incident
directions in both blocks but no mixed square. Thus the union is
cornerless. A contraction or root-containing section in either block
makes that block peel to the common root; remove it except for the root,
then peel the other block. The opposite section leaves a positive minor
of one block. Every elementary proper minor is therefore positive, proving
forbiddenness. The verifier also replays explicit orders for all $72$
minors of this double, stored in \path{double527.json} in the archive.
\end{proof}

The same certificates verify $24$, $12$, and $12$ distinct labelled rooted
inputs of orders $266$, $265$, and $264$, respectively; isomorphism types
have not been counted. Pairwise root gluing is sound for all of them.
This improves the earlier $575$-vertex cut-vertex construction, not the
unrestricted upper bound $267$. It does not establish that deleting a
forced corner and changing the root preserves rooted minimality in general.

%======================================================================
\begin{remark}[Signed coordinate identifications of the rooted input]
\label{rem:root264-identifications}
A coordinate identification retains the words satisfying
$x_i\mathbin{\oplus}x_j=b$ and then suppresses coordinate $j$.
It need not be a pc-minor operation, so the proper-minor endpoint assertion
above does not cover it. All $132$ choices $(i,j,b)$ for the explicit
$264$-vertex input were checked. Exactly $15$ outputs are ample, and every
vertex of each is attainable as a peeling endpoint, with $1941$ full orders
in \path{evidence/root264_identifications/}. The other $117$ outputs are
nonample and cannot directly supply ample rooted inputs.

In particular, the identifications $(i,j,b)=(1,10,0)$ and $(3,8,0)$ give
ample classes of orders $107$ and $108$. In fact each peels to every
contained cube. The numbers of vertices, edges, squares and three-cubes are
$(107,243,179,42)$ and $(108,246,182,43)$, respectively; there are no
higher-dimensional cubes. All $1150$ corresponding relative orders are
verified, with cube coverage independently enumerated from incident edges.

Consequently neither class can occur as a piece obtained by cutting an
ample forbidden pc-minor along a separating cube: the separator theorem of
\cite{BousquetDampReader} requires every piece to fail peeling to the common
cube. In particular, any number of coordinate-block copies glued along a
common cube peel, by clearing their private vertices and then the cube.
These checks do not cover repairs of the nonample identification outputs,
repeated identifications, or arbitrary small classes.
\end{remark}

%======================================================================
\section{Limitations and open problems}
\label{sec:open}

\paragraph{What is and is not excluded.}
Theorems~\ref{thm:A}--\ref{thm:C} exclude three kinds of small cornerless
classes: maximum ones on few coordinates, arbitrary ones on at most seven
coordinates, and one-coordinate lifts of Hall's contractions.  A cornerless
ample class of order between \(64\) and \(266\), if it exists, uses at least
eight coordinates, is not maximum on eight coordinates, and is not a lift
of a Hall contraction.  It could be a lift of a quite different base: for
any coordinate \(z\) of a cornerless ample \(Y\), the base \(\pi_zY\) is an
ample class which, by Theorem~\ref{thm:B}, lives on at least seven
coordinates.

\paragraph{Eight coordinates.}
The next case of Theorem~\ref{thm:B} is \(n=8\).  The formulas \(F_{8,d}\)
are roughly four times larger than for \(n=7\), and the hardest cases for
\(n=7\) already needed hours of solver time.  One simplification is valid
for fixed \(d\): only sets of size at most \(d+1\) need the trace
clauses~(Y), since a larger set outside \(\Sigma\) contains a
\((d+1)\)-set outside \(\Sigma\).  Splitting on the membership of a few
\(d\)-sets in \(\Sigma\) gives independent subformulas.  These computations
are in progress, and we claim no part of the case \(n=8\).

\begin{problem}
Does every nonempty ample class on eight coordinates have a corner?
\end{problem}

\paragraph{Maximum classes.}
Theorem~\ref{thm:A} leaves maximum classes on nine to eleven coordinates.
Only finitely many pairs \((d,n)\) with \(\Phi_d(n)<299\) remain, and a
maximum class smaller than Hall's would have to be one of them.

\begin{problem}
Is Hall's class a smallest cornerless maximum class?
\end{problem}

\paragraph{The interval.}
The central problem remains to close \(64\le\cmin\le267\).

\begin{problem}
Give a description of a cornerless ample class with the complex of
Proposition~\ref{prop:record-complex} that can be verified by hand.  Does the
parity construction generalize to other numbers of groups or other
codes, and can such a generalization produce cornerless ample classes of
order below \(267\)?
\end{problem}  Our results
suggest that a smaller example, if any, is not close to Hall's class in the
sense of Theorem~\ref{thm:C}, and that the lower bound should be approached
through the number of coordinates as well as the order.

%======================================================================
\section*{Acknowledgements}
Hall's class is used in the form published with
\cite{ChalopiChepoiMoran22}.

\bibliographystyle{plain}
\bibliography{main,references}

\end{document}
